\documentclass[UTF8,a4paper,11pt]{ctexart}

\usepackage{amsmath,amssymb,amsfonts,amsthm,mathtools,bm}
\usepackage{geometry}
\usepackage{booktabs,array,longtable,multirow}
\usepackage{enumitem}
\usepackage{xcolor}
\usepackage{hyperref}
\usepackage[capitalise,noabbrev]{cleveref}
\usepackage{microtype}
\usepackage{mathrsfs}

\geometry{left=2.35cm,right=2.35cm,top=2.35cm,bottom=2.35cm}
\hypersetup{
  colorlinks=true,
  linkcolor=blue!60!black,
  citecolor=blue!60!black,
  urlcolor=blue!60!black,
  pdftitle={线性代数习题课整理稿（1班，第1--13次）}
}
\setlist{itemsep=0.25em,topsep=0.4em}
\setlength{\parindent}{2em}
\setlength{\parskip}{0.2em}
\allowdisplaybreaks[3]
\numberwithin{equation}{section}

\newcommand{\F}{\mathbb F}
\newcommand{\R}{\mathbb R}
\newcommand{\C}{\mathbb C}
\newcommand{\Q}{\mathbb Q}
\newcommand{\Z}{\mathbb Z}
\newcommand{\N}{\mathbb N}
\newcommand{\Span}{\operatorname{span}}
\newcommand{\rank}{\operatorname{rank}}
\newcommand{\Ker}{\operatorname{ker}}
\newcommand{\Image}{\operatorname{Im}}
\newcommand{\tr}{\operatorname{tr}}
\newcommand{\diag}{\operatorname{diag}}
\newcommand{\sgn}{\operatorname{sgn}}
\newcommand{\adj}{\operatorname{adj}}
\newcommand{\End}{\operatorname{End}}
\newcommand{\Aut}{\operatorname{Aut}}
\newcommand{\Hom}{\operatorname{Hom}}
\newcommand{\GL}{\operatorname{GL}}
\newcommand{\Id}{I}
\newcommand{\ones}{\bm 1}
\newcommand{\dd}{\,\mathrm d}
\newcommand{\ii}{\mathrm i}
\newcommand{\e}{\mathrm e}
\newcommand{\norm}[1]{\left\lVert #1\right\rVert}
\newcommand{\abs}[1]{\left\lvert #1\right\rvert}
\newcommand{\set}[1]{\left\{#1\right\}}
\newcommand{\inner}[2]{\left\langle #1,#2\right\rangle}

\theoremstyle{plain}
\newtheorem{theorem}{定理}[section]
\newtheorem{proposition}[theorem]{命题}
\newtheorem{lemma}[theorem]{引理}
\newtheorem{corollary}[theorem]{推论}
\theoremstyle{definition}
\newtheorem{definition}[theorem]{定义}
\newtheorem{exercise}[theorem]{习题}
\newtheorem{example}[theorem]{例}
\theoremstyle{remark}
\newtheorem{remark}[theorem]{注}
\newenvironment{solution}{\begin{proof}[解]}{\end{proof}}

\title{线性代数习题课整理稿\\\large 1班，第1--13次}
\author{依据手写讲义统一排版}
\date{}

\begin{document}
\maketitle
\begin{abstract}
本文档依据十三份手写习题课讲义整理而成。原稿中的定义、命题、典型例题与主要推导均按原顺序重写，重复的初等行变换适度压缩，但保留判别方法、关键中间式与最终结论。全文采用统一符号与版式，并按每份源文件分别存放于独立的 \texttt{sections/XX.tex} 文件中，由主文件依次通过 \verb|\input| 引入。
\end{abstract}
\tableofcontents
\clearpage

\section{习题课 \#1（9月13日）}

\subsection{Gauss--Jordan 消元与线性方程组}

本次课首先讨论含参数的非齐次线性方程组。核心工具是将方程组写成增广矩阵，并通过可逆的初等行变换化为行阶梯形或简化行阶梯形。

\begin{example}
考虑
\begin{equation}
\begin{cases}
 x_1-x_2+x_3=1,\\
 x_1-x_2-x_3=3,\\
 2x_1-2x_2-x_3=5.
\end{cases}
\end{equation}
其增广矩阵为
\begin{equation}
\left(
\begin{array}{ccc|c}
1&-1&1&1\\
1&-1&-1&3\\
2&-2&-1&5
\end{array}
\right).
\end{equation}
经初等行变换可化为
\begin{equation}
\left(
\begin{array}{ccc|c}
1&-1&0&2\\
0&0&1&-1\\
0&0&0&0
\end{array}
\right).
\end{equation}
因此解的一般形式为
\begin{equation}
 x_1=x_2+2,
 \qquad x_3=-1,
 \qquad x_2\in\F.
\end{equation}
\end{example}

\begin{definition}
矩阵的三类初等行变换为
\begin{enumerate}
  \item 交换两行；
  \item 将某一行乘以非零常数；
  \item 将某一行的常数倍加到另一行。
\end{enumerate}
它们都可逆，因此不会改变线性方程组的解集。
\end{definition}

\begin{remark}
消元过程中每一步都必须是等价变形。对于不等式，这一点尤其容易被忽略。例如
\begin{equation}
 \frac{2x+3}{x+4}\leq 1
 \quad\Longleftrightarrow\quad
 \frac{x-1}{x+4}\leq 0,
\end{equation}
不能在不知道 $x+4$ 符号的情况下直接同乘 $x+4$。正确答案为
\begin{equation}
 x\in(-4,1].
\end{equation}
同理，在求方程 $f(x)=0$ 的解集时，若某一步只保留了单向蕴含，就必须额外检查增根或漏根。Gauss--Jordan 消元之所以可靠，正是因为每一步初等行变换均可逆。
\end{remark}

\subsection{解的三种可能与数域的影响}

线性方程组的解集只有三种可能：无解、唯一解或无穷多解。行化简后可按如下方式判断。
\begin{enumerate}
  \item 若出现形如 $(0,\ldots,0\mid c)$ 且 $c\neq0$ 的行，则方程组无解。
  \item 若每个未知量所在列都有主元，则方程组至多有一个解；在相容时即为唯一解。
  \item 若方程组相容且存在自由变量，则有无穷多解。
\end{enumerate}

所取数域也会影响结论。例如同一个整数系数方程组，在 $\Q$、$\R$ 与有限域 $\F_p$ 上可能具有不同的秩和不同的解数。行变换中的除法只能除以当前数域中的非零元素。

\subsection{典型计算}

\begin{exercise}[线性规划的一个简单例子]
设三种原料的某种成分含量分别为 $12\%$、$15\%$ 和 $22\%$。若总量固定，并且第三种原料的用量固定为 $5$，则约束可整理为
\begin{equation}
 x_1+x_2=5,
 \qquad x_3=5,
 \qquad x_1,x_2,x_3\geq0.
\end{equation}
目标函数为
\begin{equation}
 L=0.12x_1+0.15x_2+0.22x_3.
\end{equation}
求 $L$ 的可能范围。
\end{exercise}

\begin{solution}
由 $x_1=5-x_2$ 和 $0\leq x_2\leq5$，
\begin{equation}
 L=0.12(5-x_2)+0.15x_2+0.22\cdot5
   =1.7+0.03x_2.
\end{equation}
故
\begin{equation}
 1.7\leq L\leq1.85.
\end{equation}
\end{solution}

\begin{exercise}
设 $n\geq2$，$A$ 是元素均为 $\pm1$ 的 $n\times n$ 整数矩阵。证明 $\det A$ 为偶数。
\end{exercise}

\begin{proof}
模 $2$ 意义下有 $1\equiv-1$，故 $A$ 的所有元素均同余于 $1$。于是 $A$ 在 $\F_2$ 上的各行相同，秩至多为 $1<n$，从而
\begin{equation}
 \det A\equiv0\pmod2.
\end{equation}
故 $\det A$ 为偶数。
\end{proof}

\begin{exercise}[插值方程的相容性]
问是否存在二次多项式
\begin{equation}
 p(x)=ax^2+bx+c
\end{equation}
同时通过四个给定点。将插值条件写成线性方程组后，可用增广矩阵判断相容性。
\end{exercise}

\begin{solution}
以原稿中的点
\begin{equation}
 (0,2),\quad (-4,1),\quad (-1,3),\quad (1,2)
\end{equation}
为例，代入后得到
\begin{equation}
\begin{cases}
 c=2,\\
 16a-4b+c=1,\\
 a-b+c=3,\\
 a+b+c=2.
\end{cases}
\end{equation}
由后两式及 $c=2$ 得
\begin{equation}
 a-b=1,
 \qquad
 a+b=0,
\end{equation}
从而 $a=\tfrac12$、$b=-\tfrac12$。但此时
\begin{equation}
 16a-4b+c=8+2+2=12\neq1,
\end{equation}
故方程组不相容，不存在满足条件的二次多项式。这个例子说明，当插值条件的数目超过自由系数数目时，必须额外检查数据之间的相容性。
\end{solution}

\begin{exercise}[一类对称耦合方程组]
设 $a_i\neq0$，且
\begin{equation}
 1+\sum_{i=1}^n a_i^{-1}\neq0.
\end{equation}
求解
\begin{equation}
 (1+a_i)x_i+\sum_{j\neq i}x_j=b_i,
 \qquad i=1,\ldots,n.
\end{equation}
\end{exercise}

\begin{solution}
令
\begin{equation}
 X_0=\sum_{j=1}^n x_j.
\end{equation}
第 $i$ 个方程化为
\begin{equation}
 a_i x_i+X_0=b_i,
 \qquad
 x_i=a_i^{-1}(b_i-X_0).
\end{equation}
求和得
\begin{equation}
 X_0=\sum_{i=1}^n a_i^{-1}b_i-X_0\sum_{i=1}^n a_i^{-1},
\end{equation}
故
\begin{equation}
 X_0=\frac{\sum_{i=1}^n a_i^{-1}b_i}
 {1+\sum_{i=1}^n a_i^{-1}}.
\end{equation}
因此
\begin{equation}
 x_i=a_i^{-1}b_i-
 \frac{a_i^{-1}\sum_{j=1}^n a_j^{-1}b_j}
 {1+\sum_{j=1}^n a_j^{-1}},
 \qquad i=1,\ldots,n.
\end{equation}
\end{solution}

\begin{exercise}[含参数非齐次方程组]
讨论
\begin{equation}
\begin{cases}
 x_1+x_2+x_3+x_4+x_5=1,\\
 3x_1+2x_2+x_3+x_4-3x_5=a,\\
 x_2+2x_3+2x_4+6x_5=3,\\
 5x_1+4x_2+3x_3+3x_4-x_5=b
\end{cases}
\end{equation}
的相容性。
\end{exercise}

\begin{solution}
对增广矩阵作消元，最后两条独立的相容性条件为
\begin{equation}
 a=0,
 \qquad b=2.
\end{equation}
在此条件下有三个自由变量，可取 $x_3,x_4,x_5$，并得到
\begin{equation}
 x_1=x_3+x_4+5x_5-2,
 \qquad
 x_2=-2x_3-2x_4-6x_5+3.
\end{equation}
若 $(a,b)\neq(0,2)$，消元中出现矛盾行，方程组无解。
\end{solution}

\begin{exercise}[含参数齐次方程组]
讨论
\begin{equation}
\begin{cases}
 x_1+x_2+x_3=0,\\
 2x_1+x_2-a x_3=0,\\
 x_1-2x_2-3x_3=0
\end{cases}
\end{equation}
的解。
\end{exercise}

\begin{solution}
系数矩阵行列式为
\begin{equation}
 \det\begin{pmatrix}
 1&1&1\\
 2&1&-a\\
 1&-2&-3
 \end{pmatrix}
 =-3a-2.
\end{equation}
因此当 $a\neq-\frac23$ 时只有零解；当 $a=-\frac23$ 时秩下降，存在一维非零解空间。
\end{solution}

\begin{exercise}[两个行列式恒等式]
设 $a_i,b_i,c_i\in\F$。证明
\begin{equation}
 \det\begin{pmatrix}
 a_1-b_1&b_1-c_1&c_1-a_1\\
 a_2-b_2&b_2-c_2&c_2-a_2\\
 a_3-b_3&b_3-c_3&c_3-a_3
 \end{pmatrix}=0,
\end{equation}
并证明
\begin{equation}
 \det\begin{pmatrix}
 a_1+b_1&b_1+c_1&c_1+a_1\\
 a_2+b_2&b_2+c_2&c_2+a_2\\
 a_3+b_3&b_3+c_3&c_3+a_3
 \end{pmatrix}
 =2\det\begin{pmatrix}
 a_1&b_1&c_1\\
 a_2&b_2&c_2\\
 a_3&b_3&c_3
 \end{pmatrix}.
\end{equation}
\end{exercise}

\begin{proof}
第一个矩阵三列之和为零，故列向量线性相关。对第二个矩阵，记右端三列为 $a,b,c$，则左端三列为 $a+b,b+c,c+a$。相应的列变换矩阵为
\begin{equation}
 T=\begin{pmatrix}1&0&1\\1&1&0\\0&1&1\end{pmatrix},
 \qquad \det T=2,
\end{equation}
所以左端行列式等于 $2\det(a,b,c)$。
\end{proof}

\begin{exercise}[由两个已知解反推参数]
已知
\begin{equation}
 (0,1,0)^T,
 \qquad (-3,2,2)^T
\end{equation}
均为方程组
\begin{equation}
\begin{cases}
 x_1-x_2+2x_3=-1,\\
 3x_1+x_2+4x_3=1,\\
 ax_1+bx_2+cx_3=d
\end{cases}
\end{equation}
的解。求使该方程组有无穷多解的参数条件。
\end{exercise}

\begin{solution}
代入第一个解得 $b=d$；代入第二个解得
\begin{equation}
 -3a+2b+2c=d.
\end{equation}
结合 $b=d$ 可得
\begin{equation}
 3a=b+2c.
\end{equation}
前两行系数向量线性无关，因此它们给出秩为 $2$ 的相容系统。第三行要使秩不增加，必须是前两行的线性组合，这恰好等价于
\begin{equation}
 b=d,
 \qquad 3a=b+2c.
\end{equation}
在此条件下未知数个数为 $3$ 而秩为 $2$，故解集为一条仿射直线。
\end{solution}

\section{习题课 \#2（9月20日）}

\subsection{余子式、伴随矩阵与 Cramer 法则}

\begin{definition}
设 $A=(a_{ij})\in\F^{n\times n}$。删去第 $i$ 行和第 $j$ 列所得子式记为
\begin{equation}
 M_{ij},
\end{equation}
相应的代数余子式为
\begin{equation}
 C_{ij}=(-1)^{i+j}M_{ij}.
\end{equation}
由 $C_{ij}$ 组成的矩阵 $C=(C_{ij})$ 的转置称为 $A$ 的伴随矩阵，记为
\begin{equation}
 \adj(A)=C^T.
\end{equation}
\end{definition}

\begin{theorem}
对任意 $1\leq i,k\leq n$，有
\begin{equation}
 \sum_{j=1}^n a_{kj}C_{ij}=\delta_{ik}\det A.
\end{equation}
因此
\begin{equation}
 A\adj(A)=\adj(A)A=(\det A)I_n.
\end{equation}
若 $\det A\neq0$，则
\begin{equation}
 A^{-1}=\frac{1}{\det A}\adj(A).
\end{equation}
\end{theorem}

\begin{corollary}[Cramer 法则]
若 $A\in\F^{n\times n}$ 且 $\det A\neq0$，则方程组 $Ax=b$ 有唯一解，并且
\begin{equation}
 x_j=\frac{\det A_j(b)}{\det A},
 \qquad j=1,\ldots,n,
\end{equation}
其中 $A_j(b)$ 表示将 $A$ 的第 $j$ 列替换为 $b$ 后所得矩阵。
\end{corollary}

\subsection{两类 Toeplitz 行列式}

\begin{exercise}[三对角 Toeplitz 行列式]
计算
\begin{equation}
 D_n=\begin{vmatrix}
 a&b&&&\\
 c&a&b&&\\
 &c&a&\ddots&\\
 &&\ddots&\ddots&b\\
 &&&c&a
 \end{vmatrix}_{n\times n}.
\end{equation}
\end{exercise}

\begin{solution}
沿第一行展开得到递推式
\begin{equation}
 D_n=aD_{n-1}-bcD_{n-2},
 \qquad D_0=1,
 \qquad D_1=a.
\end{equation}
设 $\eta_1,\eta_2$ 为
\begin{equation}
 t^2-at+bc=0
\end{equation}
的两个根。若 $\eta_1\neq\eta_2$，则
\begin{equation}
 D_n=\frac{\eta_1^{n+1}-\eta_2^{n+1}}{\eta_1-\eta_2}.
\end{equation}
若 $\eta_1=\eta_2=\eta$，则取极限得
\begin{equation}
 D_n=(n+1)\eta^n.
\end{equation}
\end{solution}

\begin{exercise}[常上、下三角 Toeplitz 行列式]
计算
\begin{equation}
 D_n=\det\begin{pmatrix}
 a&b&b&\cdots&b\\
 c&a&b&\cdots&b\\
 c&c&a&\cdots&b\\
 \vdots&\vdots&\vdots&\ddots&\vdots\\
 c&c&c&\cdots&a
 \end{pmatrix}.
\end{equation}
\end{exercise}

\begin{solution}
将相邻两行作差，或分别沿首行、末行处理，可得到
\begin{equation}
 D_n=(a-b)D_{n-1}+b(a-c)^{n-1}
\end{equation}
以及对称的递推式
\begin{equation}
 D_n=(a-c)D_{n-1}+c(a-b)^{n-1}.
\end{equation}
当 $b\neq c$ 时，解为
\begin{equation}
 D_n=\frac{b(a-c)^n-c(a-b)^n}{b-c}.
\end{equation}
当 $b=c$ 时，由极限或直接递推可得
\begin{equation}
 D_n=(a-b)^{n-1}\bigl(a+(n-1)b\bigr).
\end{equation}
\end{solution}

\subsection{Cauchy 行列式与循环矩阵}

\begin{theorem}[Cauchy 行列式]
设 $x_i+y_j\neq0$，且 $x_i$ 两两不同、$y_j$ 两两不同，则
\begin{equation}
 \det\left(\frac{1}{x_i+y_j}\right)_{i,j=1}^n
 =
 \frac{
 \displaystyle\prod_{1\leq i<j\leq n}(x_j-x_i)(y_j-y_i)}
 {\displaystyle\prod_{i=1}^n\prod_{j=1}^n(x_i+y_j)}.
\end{equation}
\end{theorem}

\begin{proof}
将第 $j$ 列与第一列作适当差分，可逐步提出因子 $y_j-y_1$；再对行作同样处理，可提出 $x_j-x_1$。余下行列式降为 $n-1$ 阶 Cauchy 行列式，归纳即得结论。另一种写法是将分母统一后观察分子关于 $x_i,y_j$ 分别为交错多项式，因此必含两个 Vandermonde 因子，再比较最高次项常数。
\end{proof}

\begin{theorem}[循环矩阵的特征值]
令
\begin{equation}
 C=\begin{pmatrix}
 a_0&a_1&\cdots&a_{n-1}\\
 a_{n-1}&a_0&\cdots&a_{n-2}\\
 \vdots&\vdots&\ddots&\vdots\\
 a_1&a_2&\cdots&a_0
 \end{pmatrix},
 \qquad
 f(z)=\sum_{j=0}^{n-1}a_jz^j,
\end{equation}
并取 $\omega=\e^{2\pi\ii/n}$。则
\begin{equation}
 v_k=(1,\omega^k,\omega^{2k},\ldots,\omega^{(n-1)k})^T
\end{equation}
是特征向量，对应特征值 $f(\omega^k)$。因此
\begin{equation}
 \det C=\prod_{k=0}^{n-1}f(\omega^k).
\end{equation}
\end{theorem}

\begin{exercise}[正弦行列式]
计算
\begin{equation}
 D=\det\bigl(\sin(j\theta_i)\bigr)_{i,j=1}^n.
\end{equation}
\end{exercise}

\begin{solution}
利用 Chebyshev 多项式第二类
\begin{equation}
 \sin(j\theta)=\sin\theta\,U_{j-1}(\cos\theta),
\end{equation}
且 $U_{j-1}$ 的首项系数为 $2^{j-1}$，于是
\begin{equation}
 D
 =2^{n(n-1)/2}
 \left(\prod_{i=1}^n\sin\theta_i\right)
 \left(\prod_{1\leq i<j\leq n}(\cos\theta_j-\cos\theta_i)\right).
\end{equation}
这与原稿中利用复指数和 Vandermonde 行列式的推导等价。
\end{solution}

\subsection{行列式的求导与几个加分题}

\begin{theorem}[行列式求导公式]
设 $A(t)=(a_{ij}(t))$ 的各元素可微。令 $B_j(t)$ 表示将 $A(t)$ 的第 $j$ 列替换为其导数列后所得矩阵，则
\begin{equation}
 \frac{\dd}{\dd t}\det A(t)
 =\sum_{j=1}^n\det B_j(t)
 =\sum_{i,j=1}^n a'_{ij}(t)C_{ij}(t).
\end{equation}
若 $A(t)$ 可逆，还可写为 Jacobi 公式
\begin{equation}
 \frac{\dd}{\dd t}\det A(t)
 =\det A(t)\,\tr\bigl(A(t)^{-1}A'(t)\bigr).
\end{equation}
\end{theorem}

\begin{proof}
行列式对每一列均为线性函数。对乘积形式逐列求导，恰好得到把某一列替换为导数列的 $n$ 项之和。再对每个 $\det B_j$ 按第 $j$ 列展开，即得代数余子式表达式。
\end{proof}

\begin{exercise}[带参数的 Vandermonde 行列式]
令
\begin{equation}
 D_1(y)=
 \det\begin{pmatrix}
 1&1&\cdots&1\\
 y&x_1&\cdots&x_n\\
 y^2&x_1^2&\cdots&x_n^2\\
 \vdots&\vdots&&\vdots\\
 y^n&x_1^n&\cdots&x_n^n
 \end{pmatrix}.
\end{equation}
求其关于 $y$ 的展开。
\end{exercise}

\begin{solution}
Vandermonde 公式给出
\begin{equation}
 D_1(y)=
 \prod_{j=1}^n(x_j-y)
 \prod_{1\leq i<j\leq n}(x_j-x_i).
\end{equation}
若 $e_r(x_1,\ldots,x_n)$ 表示第 $r$ 个初等对称多项式，则
\begin{equation}
 \prod_{j=1}^n(x_j-y)
 =\sum_{k=0}^n(-1)^k e_{n-k}(x_1,\ldots,x_n)y^k.
\end{equation}
因此 $y^k$ 的系数为
\begin{equation}
 (-1)^k e_{n-k}(x_1,\ldots,x_n)
 \prod_{i<j}(x_j-x_i).
\end{equation}
\end{solution}

\begin{exercise}
设
\begin{equation}
 D_2=\det\bigl(1+x_i^j\bigr)_{i,j=1}^n.
\end{equation}
证明
\begin{equation}
 D_2=
 \left(2\prod_{k=1}^n x_k-\prod_{k=1}^n(x_k-1)\right)
 \prod_{1\leq i<j\leq n}(x_j-x_i).
\end{equation}
\end{exercise}

\begin{proof}
在行列式左侧添一行一列，将其写成
\begin{equation}
 D_2=
 \det\begin{pmatrix}
 1&1&1&\cdots&1\\
 -1&x_1&x_1^2&\cdots&x_1^n\\
 \vdots&\vdots&\vdots&&\vdots\\
 -1&x_n&x_n^2&\cdots&x_n^n
 \end{pmatrix}.
\end{equation}
把第一列拆成 $(-1,1,\ldots,1)^T$ 与一个只在首项非零的向量之差，并分别使用 Vandermonde 公式，即得所述表达式。
\end{proof}

\begin{exercise}[全体元素同时加参数]
设 $A=(a_{ij})$，$J=\ones\ones^T$，证明
\begin{equation}
 \det(A+tJ)=\det A+t\sum_{i,j=1}^n C_{ij}(A).
\end{equation}
\end{exercise}

\begin{proof}
矩阵 $tJ$ 的各列相同。由行列式对列的多重线性，展开 $\det(A+tJ)$ 时，含有两列以上来自 $tJ$ 的项均为零，因此只剩常数项和一次项。一次项正是所有代数余子式之和。也可直接使用矩阵行列式引理。
\end{proof}

\begin{exercise}[参数特征方程]
求齐次方程组
\begin{equation}
\begin{cases}
 (\lambda-2)x_1-3x_2-2x_3=0,\\
 -x_1+(\lambda-8)x_2-2x_3=0,\\
 2x_1+14x_2+(\lambda+3)x_3=0
\end{cases}
\end{equation}
存在非零解时的 $\lambda$。
\end{exercise}

\begin{solution}
系数行列式为
\begin{equation}
 \det\begin{pmatrix}
 \lambda-2&-3&-2\\
 -1&\lambda-8&-2\\
 2&14&\lambda+3
 \end{pmatrix}
 = (\lambda-1)(\lambda-3)^2.
\end{equation}
故且仅当
\begin{equation}
 \lambda=1\quad\text{或}\quad\lambda=3
\end{equation}
时存在非零解。
\end{solution}

\begin{remark}
一般的参数齐次方程组
\begin{equation}
 A x=\lambda x
\end{equation}
存在非零解的参数由特征多项式 $\det(\lambda I-A)$ 的零点给出。若该多项式不恒为零，则不同参数值至多有 $n$ 个；在代数闭域中按重数恰有 $n$ 个。
\end{remark}

\subsection{分段插值与严格对角占优型行列式}

手稿最后补充了分段多项式插值。若在 $n$ 个区间上分别使用次数不超过 $s$ 的多项式，则未知系数共有 $n(s+1)$ 个。插值端点条件、区间连接处的连续性条件以及边界条件必须在数量与独立性上共同保证唯一性。分段线性插值对应 $s=1$ 且只要求函数连续；三次样条对应 $s=3$、连接处要求到二阶导数连续，还需补充两个边界条件，例如自然边界条件。

\begin{theorem}[一类 $M$-矩阵判别]
设实矩阵 $A=(a_{ij})\in\R^{n\times n}$ 满足
\begin{equation}
 a_{ii}>0,
 \qquad a_{ij}<0\quad(i\neq j),
 \qquad \sum_{j=1}^n a_{ij}>0\quad(i=1,\ldots,n).
\end{equation}
则
\begin{equation}
 \det A>0.
\end{equation}
\end{theorem}

\begin{proof}
对 $n$ 归纳。以 $a_{11}>0$ 为主元消去第一列。Schur 补的非对角元仍为负，而其每一行的元素和为
\begin{equation}
 \sum_{j=2}^n\left(a_{ij}-a_{i1}a_{11}^{-1}a_{1j}\right)
 =\sum_{j=1}^n a_{ij}
 -a_{i1}a_{11}^{-1}\sum_{j=1}^n a_{1j}>0,
\end{equation}
因为 $a_{i1}<0$ 且第一行行和为正。故 Schur 补满足同样条件，由归纳假设其行列式为正。于是
\begin{equation}
 \det A=a_{11}\det(A/a_{11})>0.
\end{equation}
\end{proof}

\section{习题课 \#3（9月27日）}

\subsection{线性相关、线性无关与交换定理}

\begin{remark}
当 $n\geq2$ 时，$S_n$ 中奇排列与偶排列数目相同。一个简洁证明是
\begin{equation}
 \sum_{\sigma\in S_n}\sgn(\sigma)
 =\det\begin{pmatrix}
 1&\cdots&1\\
 \vdots&&\vdots\\
 1&\cdots&1
 \end{pmatrix}=0.
\end{equation}
\end{remark}

\begin{theorem}[方阵列向量线性无关的等价条件]
设 $A=(\alpha_1,\ldots,\alpha_n)\in\F^{n\times n}$。下列条件等价：
\begin{enumerate}
  \item $\alpha_1,\ldots,\alpha_n$ 线性无关；
  \item $\sum_{j=1}^n c_j\alpha_j=0$ 只有零解；
  \item $\det A\neq0$；
  \item 对任意 $b\in\F^n$，方程 $Ax=b$ 有唯一解；
  \item $A$ 可逆。
\end{enumerate}
\end{theorem}

\begin{remark}
对于 $s\leq n$ 个向量 $\alpha_1,\ldots,\alpha_s\in\F^n$，线性无关仍等价于齐次关系只有零解，但不再能直接使用一个 $s\times s$ 行列式，除非选取适当的 $s$ 行形成非零子式。
\end{remark}

\begin{proposition}[用线性组合替换一个向量]
设 $\alpha_1,\ldots,\alpha_s$ 线性无关，并令
\begin{equation}
 \beta=\sum_{i=1}^s k_i\alpha_i.
\end{equation}
若 $k_j\neq0$，则
\begin{equation}
 \alpha_1,\ldots,\alpha_{j-1},\beta,
 \alpha_{j+1},\ldots,\alpha_s
\end{equation}
仍线性无关。
\end{proposition}

\begin{proof}
若
\begin{equation}
 r\beta+\sum_{i\neq j}r_i\alpha_i=0,
\end{equation}
代入 $\beta$ 后，$\alpha_j$ 的系数为 $rk_j$。由原向量组线性无关且 $k_j\neq0$，先得 $r=0$，再得所有 $r_i=0$。
\end{proof}

\begin{theorem}[Steinitz 交换定理]
设 $\alpha_1,\ldots,\alpha_s$ 张成向量空间 $V$，而 $\beta_1,\ldots,\beta_t$ 线性无关。则 $t\leq s$，并且可以从 $\alpha_1,\ldots,\alpha_s$ 中删去 $t$ 个向量，使
\begin{equation}
 \beta_1,\ldots,\beta_t
\end{equation}
与余下的 $s-t$ 个 $\alpha$ 向量仍张成 $V$。若原来的 $\alpha$ 向量组是一组基，则替换后仍是一组基。
\end{theorem}

\begin{proof}
对 $t$ 归纳。因 $\beta_1\in\Span(\alpha_1,\ldots,\alpha_s)$，其表示中至少有一个系数非零。由上一命题可用 $\beta_1$ 替换相应的 $\alpha$ 而不改变张成空间。随后将 $\beta_2$ 写成新生成组的线性组合。若 $\beta_2$ 在 $\beta_1$ 之外的所有系数均为零，则 $\beta_1,\beta_2$ 相关，矛盾；故还可替换另一个 $\alpha$。如此继续即可。
\end{proof}

\begin{remark}
有限向量组可按其张成空间比较：
\begin{equation}
 A\preccurlyeq B
 \quad\Longleftrightarrow\quad
 \Span A\subseteq\Span B.
\end{equation}
这是一个预序。将张成同一子空间的向量组视为等价后，所得偏序正是有限生成子空间按包含关系构成的偏序。
\end{remark}

\subsection{秩与线性关系的支撑}

\begin{proposition}
设 $\alpha_1,\ldots,\alpha_s\in\F^n$，并固定 $r\leq\min\{n,s\}$。下列条件等价：
\begin{enumerate}
  \item 任取其中 $r$ 个向量均线性无关；
  \item 每个非平凡线性关系
  \begin{equation}
   \sum_{j=1}^s x_j\alpha_j=0
  \end{equation}
  至少含有 $r+1$ 个非零系数。
\end{enumerate}
\end{proposition}

\begin{proof}
若存在只含不超过 $r$ 个非零系数的非平凡关系，则把相应向量补足到 $r$ 个便得到相关组。反之，若某 $r$ 个向量相关，则其非平凡关系扩充为整个向量组上的关系，非零系数不超过 $r$。
\end{proof}

\begin{proposition}[严格对角占优判别]
设 $A=(a_{ij})\in\F^{n\times n}$，并假设绝对值来自 $\R$ 或 $\C$。若
\begin{equation}
 |a_{ii}|>\sum_{j\neq i}|a_{ij}|,
 \qquad i=1,\ldots,n,
\end{equation}
则 $A$ 可逆。
\end{proposition}

\begin{proof}
若 $Ax=0$ 且 $x\neq0$，取 $k$ 使 $|x_k|=\max_j|x_j|>0$。第 $k$ 行给出
\begin{equation}
 |a_{kk}|\,|x_k|
 \leq\sum_{j\neq k}|a_{kj}|\,|x_j|
 \leq\left(\sum_{j\neq k}|a_{kj}|\right)|x_k|,
\end{equation}
与严格不等式矛盾。
\end{proof}

\subsection{若干行列式恒等式}

\begin{exercise}[行和与列差]
设 $A=(a_{ij})\in\F^{n\times n}$，并令
\begin{equation}
 s_i=\sum_{k=1}^n a_{ik},
 \qquad
 B=(s_i-a_{ij})_{i,j=1}^n.
\end{equation}
证明
\begin{equation}
 \det B=(-1)^n(1-n)\det A.
\end{equation}
\end{exercise}

\begin{proof}
记 $J=\ones\ones^T$。因为 $A\ones=(s_1,\ldots,s_n)^T$，故
\begin{equation}
 B=A(J-I_n).
\end{equation}
矩阵 $J-I_n$ 在 $\Span\{\ones\}$ 上的特征值为 $n-1$，在其余 $n-1$ 维子空间上的特征值为 $-1$。于是
\begin{equation}
 \det(J-I_n)=(n-1)(-1)^{n-1}=(-1)^n(1-n),
\end{equation}
结论随即成立。
\end{proof}

\begin{exercise}[列差行列式与余子式之和]
设 $A=(a_{ij})\in\F^{n\times n}$，列向量为 $a_1,\ldots,a_n$。证明
\begin{equation}
 \det(a_1-a_2,a_2-a_3,\ldots,a_{n-1}-a_n,\ones)
 =\sum_{i,j=1}^n C_{ij}(A).
\end{equation}
\end{exercise}

\begin{proof}
由矩阵行列式引理的多项式形式，
\begin{equation}
 \det(A+t\ones\ones^T)
 =\det A+t\ones^T\adj(A)\ones
 =\det A+t\sum_{i,j}C_{ij}(A).
\end{equation}
另一方面，对列作连续差分可把 $A+t\ones\ones^T$ 的一次项系数化为题设左端行列式，故二者相等。也可直接反复使用列的线性性展开。
\end{proof}

\begin{exercise}[正弦和矩阵的秩]
计算
\begin{equation}
 D_n=\det\bigl(\sin(\alpha_k+\beta_\ell)\bigr)_{k,\ell=1}^n.
\end{equation}
\end{exercise}

\begin{solution}
由加法公式
\begin{equation}
 \sin(\alpha_k+\beta_\ell)
 =\sin\alpha_k\cos\beta_\ell+
   \cos\alpha_k\sin\beta_\ell,
\end{equation}
整个矩阵是两个秩一矩阵之和，故秩至多为 $2$。因此
\begin{equation}
 D_n=0,
 \qquad n\geq3.
\end{equation}
低阶情形为
\begin{equation}
 D_1=\sin(\alpha_1+\beta_1),
\end{equation}
以及
\begin{equation}
 D_2
 =\sin(\alpha_1-\alpha_2)
  \sin(\beta_2-\beta_1).
\end{equation}
原稿也用复指数展开说明：每一列均落在由
$(\e^{\ii\alpha_k})_k$ 与 $(\e^{-\ii\alpha_k})_k$ 张成的二维空间中。
\end{solution}

\begin{exercise}[$\mu$-循环行列式]
设
\begin{equation}
 A=\begin{pmatrix}
 c_0&c_1&\cdots&c_{n-1}\\
 \mu c_{n-1}&c_0&\cdots&c_{n-2}\\
 \vdots&\vdots&\ddots&\vdots\\
 \mu c_1&\mu c_2&\cdots&c_0
 \end{pmatrix},
 \qquad
 f(z)=\sum_{j=0}^{n-1}c_jz^j.
\end{equation}
计算 $\det A$。
\end{exercise}

\begin{solution}
若 $\xi^n=\mu$，令
\begin{equation}
 v(\xi)=(1,\xi,\ldots,\xi^{n-1})^T.
\end{equation}
直接计算得
\begin{equation}
 Av(\xi)=f(\xi)v(\xi).
\end{equation}
当 $\mu\neq0$ 时，方程 $\xi^n=\mu$ 的 $n$ 个根两两不同，对应的 Vandermonde 特征向量构成一组基，因此
\begin{equation}
 \det A=\prod_{\xi^n=\mu}f(\xi).
\end{equation}
该恒等式两端都是 $\mu,c_0,\ldots,c_{n-1}$ 的多项式，故也可由连续性推广到退化情形。
\end{solution}

\section{习题课 \#4（10月4日）}

\subsection{由张成空间得到的预序与秩函数}

记 $\mathcal P_{\mathrm{fin}}(\F^n)$ 为 $\F^n$ 的有限子集全体。对有限集 $A$ 定义
\begin{equation}
 \Span A=\left\{\sum_{k=1}^r c_k a_k:
 r<\infty,
 \ c_k\in\F,
 \ a_k\in A\right\}.
\end{equation}

\begin{lemma}
有限集 $A$ 中每个向量均可由 $B$ 线性表出，当且仅当
\begin{equation}
 \Span A\subseteq\Span B.
\end{equation}
\end{lemma}

因此可定义
\begin{equation}
 A\preccurlyeq B
 \quad\Longleftrightarrow\quad
 \Span A\subseteq\Span B.
\end{equation}
该关系自反且传递，但一般不反对称，所以只是预序。再定义
\begin{equation}
 A\sim B
 \quad\Longleftrightarrow\quad
 \Span A=\Span B,
\end{equation}
则商集上的 $\preccurlyeq$ 成为偏序。

\begin{definition}
有限向量组 $A$ 的秩定义为其中线性无关子组的最大基数：
\begin{equation}
 \rank A=\
 \max\{\abs{S}:S\subseteq A,\ S\text{ 线性无关}\}.
\end{equation}
\end{definition}

\begin{proposition}
若 $A\preccurlyeq B$，则
\begin{equation}
 \rank A\leq\rank B.
\end{equation}
若 $A\sim B$，则 $\rank A=\rank B$。在 $\F^n$ 中，张成整个空间的向量组秩为 $n$。
\end{proposition}

\begin{remark}
不能把 $A\preccurlyeq B$ 直接推出 $|A|\leq|B|$，因为一个向量组可含大量冗余向量。把基数替换为秩后，单调性才成立。
\end{remark}

\subsection{矩阵的行秩与列秩}

\begin{theorem}
任意矩阵的行秩等于列秩。二者共同记为 $\rank A$。
\end{theorem}

\begin{proof}
初等行变换不改变行空间的维数，同时由左乘可逆矩阵实现，因而不改变列向量之间的线性关系。将 $A$ 化为行阶梯形后，非零行数既等于行秩，也等于主元列数，即列秩。
\end{proof}

\begin{proposition}[由解集包含推出行向量包含]
设 $B\in\F^{r\times n}$，$c^T\in\F^{1\times n}$。若每个满足 $Bx=0$ 的向量也满足 $c^Tx=0$，即
\begin{equation}
 \Ker B\subseteq\Ker c^T,
\end{equation}
则 $c^T$ 是 $B$ 的行向量的线性组合。
\end{proposition}

\begin{proof}
把 $c^T$ 添在 $B$ 的末行。假设添行后秩增加，则齐次系统的零空间维数下降，从而存在 $x\in\Ker B$ 但 $c^Tx\neq0$，矛盾。因此添行不增秩，故 $c^T$ 属于 $B$ 的行空间。
\end{proof}

\subsection{向量组秩的基本估计}

设
\begin{equation}
 A=(\alpha_1,\ldots,\alpha_s),
 \qquad
 B=(\beta_1,\ldots,\beta_t).
\end{equation}
则有
\begin{align}
 \rank A&\leq n,\\
 \rank B&\leq n,\\
 \rank(A,B)&\leq\rank A+\rank B,\\
 \rank(A,\beta)&\leq\rank A+1.
\end{align}
并且
\begin{equation}
 \rank(A,B)=\rank A
 \quad\Longleftrightarrow\quad
 \beta_j\in\Span A\quad(j=1,\ldots,t).
\end{equation}

\begin{exercise}[参数矩阵的秩与极大无关列组]
讨论
\begin{equation}
 A(\lambda)=
 \begin{pmatrix}
 1&\lambda&-1&2\\
 2&-1&\lambda&5\\
 1&10&-6&1
 \end{pmatrix}
\end{equation}
的秩，并给出一组极大线性无关列。
\end{exercise}

\begin{solution}
以第一行为主元行，作
\begin{equation}
 R_2\leftarrow R_2-2R_1,
 \qquad
 R_3\leftarrow R_3-R_1.
\end{equation}
继续消元后，控制秩下降的因子为
\begin{equation}
 (\lambda+5)(\lambda-3).
\end{equation}
于是
\begin{enumerate}
  \item 当 $\lambda=3$ 时，$\rank A=2$，第 $1,2$ 列构成极大无关列组；
  \item 当 $\lambda=-5$ 时，$\rank A=3$，可取第 $1,2,4$ 列；
  \item 当 $\lambda\neq3,-5$ 时，$\rank A=3$，可取第 $1,2,3$ 列。
\end{enumerate}
特别地，$\lambda=10$ 虽使中间某个消元分母为零，但直接换主元后仍属第三种情形。
\end{solution}

\begin{proposition}[截取若干行或列后的秩下界]
设 $A\in\F^{m\times n}$ 且 $\rank A=r$。令 $A_s$ 为 $A$ 的前 $s$ 行组成的子矩阵，则
\begin{equation}
 \rank A_s\geq r+s-m.
\end{equation}
类似地，若 $B_s$ 为前 $s$ 列组成的子矩阵，则
\begin{equation}
 \rank B_s\geq r+s-n.
\end{equation}
\end{proposition}

\begin{proof}
其余 $m-s$ 行每加入一行至多使秩增加 $1$，故
\begin{equation}
 r=\rank A\leq\rank A_s+(m-s).
\end{equation}
移项即得。列的情形完全相同。
\end{proof}

\begin{example}[数值矩阵的消元]
原稿给出一个 $4\times5$ 的整数矩阵，并逐行消元。化为阶梯形后恰有三个主元，因此矩阵秩为 $3$，第 $1,2,3$ 列构成一个极大线性无关列组。此类题的要点不是保留每一步大整数运算，而是同时记录主元位置，避免在消元后误把新矩阵的列直接当作原矩阵的列。
\end{example}

\subsection{常上、下三角 Toeplitz 矩阵的秩}

考虑
\begin{equation}
 T_n(a,b,c)=
 \begin{pmatrix}
 a&b&\cdots&b\\
 c&a&\cdots&b\\
 \vdots&\vdots&\ddots&\vdots\\
 c&c&\cdots&a
 \end{pmatrix}.
\end{equation}
由上一讲的计算，
\begin{equation}
 \det T_n=
 \begin{cases}
 (a-b)^{n-1}\bigl(a+(n-1)b\bigr),&b=c,\\[0.4em]
 \displaystyle\frac{b(a-c)^n-c(a-b)^n}{b-c},&b\neq c.
 \end{cases}
\end{equation}
因此先由 $\det T_n$ 判断是否满秩；若行列式为零，再检查一个 $(n-1)$ 阶主子式即可区分秩 $n-1$ 与更低的情形。

原稿特别列出了以下退化情形：
\begin{enumerate}
  \item 若 $a=b=c\neq0$，则所有列相同，$\rank T_n=1$，任一列均为极大无关列组；
  \item 若 $b=c\neq0$ 且 $a=(1-n)b$，则 $\rank T_n=n-1$，前 $n-1$ 列可取为极大无关列组；
  \item 若 $a=b=c=0$，则 $\rank T_n=0$；
  \item 若 $a=c=0$ 且 $b\neq0$，则 $T_n$ 为严格上三角常数矩阵，$\rank T_n=n-1$，后 $n-1$ 列线性无关；
  \item 对称地，若 $a=b=0$ 且 $c\neq0$，则 $\rank T_n=n-1$，前 $n-1$ 列线性无关。
\end{enumerate}
其余情形可直接结合 $\det T_n$ 与相邻阶行列式判定。

\section{习题课 \#5（10月11日）}

\subsection{线性方程组的存在性与唯一性}

考虑线性方程组
\begin{equation}
 Ax=b,
 \qquad
 A\in\F^{m\times n},\quad x\in\F^n,\quad b\in\F^m.
\end{equation}

\begin{theorem}[相容性判别]
下列条件等价：
\begin{enumerate}
  \item 方程组 $Ax=b$ 有解；
  \item $b\in\Image A$；
  \item $\Span(A,b)=\Span A$；
  \item $\rank(A,b)=\rank A$。
\end{enumerate}
\end{theorem}

\begin{proof}
矩阵 $A$ 的列向量记为 $\alpha_1,\ldots,\alpha_n$。方程
\begin{equation}
 Ax=b
\end{equation}
等价于
\begin{equation}
 x_1\alpha_1+\cdots+x_n\alpha_n=b,
\end{equation}
所以有解当且仅当 $b$ 属于 $A$ 的列空间。把 $b$ 添作一列不改变列空间，等价于增广矩阵与系数矩阵秩相同。
\end{proof}

\begin{theorem}[唯一性判别]
若 $Ax=b$ 已知有解，则下列条件等价：
\begin{enumerate}
  \item $Ax=b$ 的解唯一；
  \item 齐次方程 $Ax=0$ 只有零解；
  \item $\Ker A=\{0\}$；
  \item $A$ 的列向量线性无关；
  \item $\rank A=n$。
\end{enumerate}
\end{theorem}

\begin{proof}
若 $x_1,x_2$ 都是解，则
\begin{equation}
 A(x_1-x_2)=0.
\end{equation}
因此非齐次方程的解是否唯一，恰由齐次方程是否存在非零解决定。其余等价关系来自秩--零化度公式
\begin{equation}
 \dim\Ker A=n-\rank A.
\end{equation}
\end{proof}

\begin{corollary}
设 $A\in\F^{m\times n}$。
\begin{enumerate}
  \item 若 $\rank A=m$，则对每个 $b\in\F^m$，方程 $Ax=b$ 都有解；
  \item 若 $\rank A=n$，则对每个 $b\in\F^m$，方程 $Ax=b$ 至多有一个解；
  \item 若 $A$ 同时满行秩和满列秩，则 $m=n$ 且 $A$ 可逆，对每个 $b$ 都有唯一解。
\end{enumerate}
\end{corollary}

\begin{remark}
Gauss--Jordan 消元给出同一判别的计算版本。阶梯形增广矩阵中出现
\begin{equation}
 (0,\ldots,0\mid c),\qquad c\neq0,
\end{equation}
时无解；若无此矛盾行，则方程相容。相容时，主元列数等于未知数个数则唯一，否则含自由变量并有无穷多解。
\end{remark}

\subsection{解集的仿射结构}

\begin{theorem}
假设 $Ax=b$ 相容，并固定一个特解 $x_0$。则全部解构成仿射子空间
\begin{equation}
 \set{x\in\F^n:Ax=b}=x_0+\Ker A
 =\set{x_0+z:z\in\Ker A}.
\end{equation}
\end{theorem}

\begin{proof}
若 $Ax=b$，则
\begin{equation}
 A(x-x_0)=0,
\end{equation}
故 $x-x_0\in\Ker A$。反之，若 $z\in\Ker A$，则
\begin{equation}
 A(x_0+z)=Ax_0+Az=b.
\end{equation}
\end{proof}

\begin{remark}
因此，非齐次方程的全部信息可以分成两部分。一个特解决定仿射子空间的位置，齐次方程的零空间决定其方向。相容方程组的自由参数个数为
\begin{equation}
 \dim\Ker A=n-\rank A.
\end{equation}
\end{remark}

\subsection{从主元列构造零空间的一组基}

设
\begin{equation}
 A=(\beta_1,\ldots,\beta_n),
 \qquad
 \rank A=r,
\end{equation}
并设 $\beta_1,\ldots,\beta_r$ 是一组极大线性无关列。对每个 $j=r+1,\ldots,n$，存在唯一系数 $c_{1j},\ldots,c_{rj}$ 使
\begin{equation}
 \beta_j=\sum_{i=1}^r c_{ij}\beta_i.
\end{equation}
令
\begin{equation}
 \eta_j=(-c_{1j},\ldots,-c_{rj},0,\ldots,0,
 \underset{j\text{ 位}}{1},0,\ldots,0)^T.
\end{equation}
则 $A\eta_j=0$。

\begin{proposition}
向量
\begin{equation}
 \eta_{r+1},\ldots,\eta_n
\end{equation}
构成 $\Ker A$ 的一组基。
\end{proposition}

\begin{proof}
每个 $\eta_j$ 都属于 $\Ker A$。它们在第 $r+1,\ldots,n$ 个坐标上形成单位矩阵，故线性无关。数量为 $n-r=\dim\Ker A$，因而是一组基。
\end{proof}

\begin{remark}
这正是把阶梯形方程中的自由变量逐个取为 $1$、其余自由变量取为 $0$ 所得到的基础解系。
\end{remark}

\subsection{行空间与零空间的正交关系}

在 $\F=\R$ 或 $\C$ 上使用标准内积。复数情形中转置均应替换为共轭转置。

\begin{theorem}[基本正交分解]
对 $A\in\R^{m\times n}$，有
\begin{equation}
 \Image A^T=(\Ker A)^\perp,
 \qquad
 \R^n=\Image A^T\oplus\Ker A.
\end{equation}
相应地，
\begin{equation}
 \Image A=(\Ker A^T)^\perp,
 \qquad
 \R^m=\Image A\oplus\Ker A^T.
\end{equation}
\end{theorem}

\begin{proof}
若 $y=A^Tu$ 且 $z\in\Ker A$，则
\begin{equation}
 y^Tz=u^TAz=0,
\end{equation}
所以 $\Image A^T\subseteq(\Ker A)^\perp$。两边维数均为
\begin{equation}
 \rank A=n-\dim\Ker A,
\end{equation}
故二者相等。标准内积正定，因此子空间与其正交补直和为全空间。
\end{proof}

\begin{remark}
在一般域上，形式 $x^Ty$ 未必正定，因而不能无条件写出上述直和。例如在 $\F_2^2$ 中，非零向量 $(1,1)^T$ 满足
\begin{equation}
 (1,1)(1,1)^T=0,
\end{equation}
一个子空间可能与自身有非零交。因此原稿中特别强调，正交直和结论依赖于实数或复数上的正定内积。
\end{remark}

\subsection{典型消元题}

\begin{exercise}
判断方程组
\begin{equation}
\begin{cases}
 3x_1-5x_2+2x_3+4x_4=2,\\
 7x_1-4x_2-x_3+3x_4=5,\\
 5x_1+7x_2-4x_3-6x_4=3
\end{cases}
\end{equation}
是否相容。
\end{exercise}

\begin{solution}
对增广矩阵作初等行变换，会出现系数部分全为零而常数项非零的一行，因此
\begin{equation}
 \rank(A,b)>\rank A.
\end{equation}
故方程组无解。
\end{solution}

\begin{exercise}
求解
\begin{equation}
\begin{cases}
 2x_1+5x_2-8x_3=8,\\
 4x_1+3x_2-9x_3=9,\\
 2x_1+3x_2-5x_3=7,\\
 x_1+8x_2-7x_3=12.
\end{cases}
\end{equation}
\end{exercise}

\begin{solution}
消元后有三个主元，且无矛盾行，故解唯一。回代得到
\begin{equation}
 (x_1,x_2,x_3)^T=(3,2,1)^T.
\end{equation}
直接代回四个方程也可验证。
\end{solution}

\begin{remark}[含参数方程组的统一处理]
原稿还对一个含参数的四元方程组逐项讨论。计算时应先在不除以可能为零的参数表达式的前提下化为阶梯形，再按使主元消失的参数值分类。每一类中分别比较 $\rank A$ 与 $\rank(A,b)$，相容时再由 $n-\rank A$ 确定自由变量数。该流程比在一开始直接除以参数更不易漏解。
\end{remark}

\subsection{有限个真子空间不能覆盖无限域上的有限维空间}

\begin{theorem}
设 $V$ 是无限域 $\F$ 上的有限维向量空间。若
\begin{equation}
 V=W_1\cup\cdots\cup W_s
\end{equation}
且每个 $W_i$ 都是线性子空间，则至少有一个 $W_i=V$。换言之，有限个真子空间不能覆盖 $V$。
\end{theorem}

\begin{proof}
先看二维情形。若 $W_1,\ldots,W_s$ 都是真子空间，则每个至多是一条过原点的直线。取线性无关向量 $u,v$。由于 $\F$ 无限，向量
\begin{equation}
 u+tv,\qquad t\in\F,
\end{equation}
给出无穷多条不同方向，而有限条直线至多包含有限个这样的方向，故不能覆盖平面。

一般维数可归纳证明。选取一个不包含在给定真子空间中的超平面，或将问题限制到适当二维仿射平面，即可归约到上述论证。
\end{proof}

\begin{remark}
无限域假设不可删去。有限域上的有限维空间本身只有有限多个向量，也可以写成有限个一维子空间的并。
\end{remark}

\subsection{两个秩结论}

\begin{proposition}[斜对称矩阵的秩为偶数]
设 $A\in\F^{n\times n}$ 满足
\begin{equation}
 A^T=-A,
\end{equation}
且 $\operatorname{char}\F\neq2$。则 $\rank A$ 为偶数。
\end{proposition}

\begin{proof}
设 $r=\rank A$。可选出一个非零的 $r$ 阶子式。对斜对称矩阵，还可通过同时交换行和列把一个非奇异主子矩阵取作
\begin{equation}
 A_r\in\F^{r\times r},
 \qquad A_r^T=-A_r.
\end{equation}
于是
\begin{equation}
 \det A_r=\det A_r^T=
 \det(-A_r)=(-1)^r\det A_r.
\end{equation}
由于 $\det A_r\neq0$ 且特征不为 $2$，必有 $(-1)^r=1$，故 $r$ 为偶数。
\end{proof}

\begin{proposition}[添一列时秩至多增加一]
对 $A\in\F^{m\times n}$ 与 $b\in\F^m$，有
\begin{equation}
 0\leq\rank(A,b)-\rank A\leq1.
\end{equation}
并且
\begin{equation}
 \rank(A,b)=\rank A
 \quad\Longleftrightarrow\quad
 Ax=b\text{ 有解}.
\end{equation}
\end{proposition}

\begin{proof}
添入一个列向量至多使列空间维数增加 $1$。秩不增加恰好意味着 $b$ 已在原列空间中。
\end{proof}

\section{习题课 \#6（10月18日）}

\subsection{Gram 矩阵与秩}

\begin{theorem}
设 $A\in\R^{m\times n}$。则
\begin{equation}
 \rank(A^TA)=\rank(AA^T)=\rank A.
\end{equation}
更一般地，若 $A\in\C^{m\times n}$，则
\begin{equation}
 \rank(A^*A)=\rank(AA^*)=\rank A.
\end{equation}
\end{theorem}

\begin{proof}
对实矩阵，若 $A^TAx=0$，则
\begin{equation}
 0=x^TA^TAx=\norm{Ax}_2^2,
\end{equation}
从而 $Ax=0$。反向包含显然成立，因此
\begin{equation}
 \Ker(A^TA)=\Ker A.
\end{equation}
由秩--零化度公式得到
\begin{equation}
 \rank(A^TA)=\rank A.
\end{equation}
对 $A^T$ 应用同一结论，并使用 $\rank A^T=\rank A$，得到
\begin{equation}
 \rank(AA^T)=\rank A.
\end{equation}
复数情形将 $T$ 换为共轭转置 $*$ 即可。
\end{proof}

\begin{remark}
上述结论依赖于正定内积。若只在一般域上使用转置，结论可能失败。例如在含有非零平方和为零的域中，可以有非零列向量 $u$ 满足 $u^Tu=0$，于是 $A=u$ 时 $A^TA=0$，但 $\rank A=1$。
\end{remark}

\begin{proposition}
设 $A\in\R^{m\times n}$。则
\begin{equation}
 \rank(AA^T)=\rank(A^TAA^T)=\rank(A^TAA^TA^T)=\cdots=\rank A.
\end{equation}
复数情形同样成立，只需把转置换为共轭转置。
\end{proposition}

\begin{proof}
每一步都可写成某个矩阵 $C$ 的 $CC^T$ 或 $C^TC$，并应用上一结论。
\end{proof}

\subsection{矩阵乘积秩等号的判别}

设
\begin{equation}
 A\in\F^{m\times n},
 \qquad
 B\in\F^{n\times p}.
\end{equation}
总有
\begin{equation}
 \rank(AB)\leq\min\{\rank A,\rank B\}.
\end{equation}
原稿进一步给出了等号成立的精确条件。

\begin{theorem}
有
\begin{equation}
 \rank(AB)=\rank A
 \quad\Longleftrightarrow\quad
 \Ker A+\Image B=\F^n.
\end{equation}
\end{theorem}

\begin{proof}
由于
\begin{equation}
 \Image(AB)=A(\Image B)\subseteq\Image A,
\end{equation}
等号成立当且仅当每个 $Ay$ 都可写成 $ABz$。这等价于对每个 $y\in\F^n$，存在 $z$ 使
\begin{equation}
 A(y-Bz)=0,
\end{equation}
也就是
\begin{equation}
 y\in\Ker A+\Image B.
\end{equation}
\end{proof}

\begin{theorem}
有
\begin{equation}
 \rank(AB)=\rank B
 \quad\Longleftrightarrow\quad
 \Ker A\cap\Image B=\{0\}.
\end{equation}
\end{theorem}

\begin{proof}
把 $A$ 限制在 $\Image B$ 上。映射
\begin{equation}
 A|_{\Image B}:\Image B\longrightarrow\Image(AB)
\end{equation}
是满射，其核为
\begin{equation}
 \Ker A\cap\Image B.
\end{equation}
由秩--零化度公式，
\begin{equation}
 \rank B=\rank(AB)+\dim(\Ker A\cap\Image B).
\end{equation}
结论随即得到。
\end{proof}

\begin{corollary}
设矩阵乘法的尺寸相容。
\begin{enumerate}
  \item 若 $\rank(AB)=\rank A$，则对任意 $C$，
  \begin{equation}
   \rank(ABC)=\rank(AC)
  \end{equation}
  并不自动成立，但若 $\Image C+\Ker(AB)$ 与相应空间满足上一判别，则可逐层判断。
  \item 若 $\rank(AB)=\rank B$，则 $A$ 在 $\Image B$ 上单射。因而对任意 $C$，
  \begin{equation}
   \rank(ABC)=\rank(BC).
  \end{equation}
  \item 若 $\rank(AB)=\rank A=\rank B$，则 $A$ 在 $\Image B$ 上给出同构
  \begin{equation}
   \Image B\cong\Image A.
  \end{equation}
\end{enumerate}
\end{corollary}

\begin{remark}
两个判别式分别描述乘法中可能损失秩的两种机制。右乘 $B$ 后若不能覆盖 $A$ 所需要的全部输入方向，则无法达到 $\rank A$；左乘 $A$ 后若 $\Image B$ 中存在落入 $\Ker A$ 的非零方向，则会损失 $\rank B$。
\end{remark}

\subsection{幂的像空间与零空间}

设 $A\in\F^{n\times n}$。显然有降链
\begin{equation}
 \F^n\supseteq\Image A\supseteq\Image A^2\supseteq\cdots
\end{equation}
和升链
\begin{equation}
 \{0\}\subseteq\Ker A\subseteq\Ker A^2\subseteq\cdots.
\end{equation}

\begin{theorem}[稳定性]
若对某个 $k\geq0$ 有
\begin{equation}
 \Image A^k=\Image A^{k+1},
\end{equation}
则对所有 $j\geq k$，
\begin{equation}
 \Image A^j=\Image A^k.
\end{equation}
同样地，若
\begin{equation}
 \Ker A^k=\Ker A^{k+1},
\end{equation}
则对所有 $j\geq k$，
\begin{equation}
 \Ker A^j=\Ker A^k.
\end{equation}
\end{theorem}

\begin{proof}
若 $\Image A^k=\Image A^{k+1}$，在两边作用 $A$ 得
\begin{equation}
 \Image A^{k+1}=\Image A^{k+2},
\end{equation}
迭代即可。

对零空间，设 $x\in\Ker A^{k+2}$。则 $Ax\in\Ker A^{k+1}=\Ker A^k$，所以
\begin{equation}
 A^{k+1}x=0,
\end{equation}
即 $x\in\Ker A^{k+1}$。反向包含本来就成立，故下一层也稳定，继而归纳。
\end{proof}

\begin{theorem}
对每个 $k\geq0$，下列条件等价：
\begin{enumerate}
  \item $\Image A^k=\Image A^{k+1}$；
  \item $\Ker A^k=\Ker A^{k+1}$；
  \item $\rank A^k=\rank A^{k+1}$。
\end{enumerate}
\end{theorem}

\begin{proof}
像空间有包含关系，维数相等即集合相等；零空间同理。又由
\begin{equation}
 \dim\Ker A^k=n-\rank A^k
\end{equation}
可知两种维数稳定等价。
\end{proof}

\begin{definition}
最小的非负整数 $k$，使
\begin{equation}
 \rank A^k=\rank A^{k+1},
\end{equation}
称为 $A$ 的指数，常记为 $\operatorname{ind}(A)$。有限维性保证该整数存在且不超过 $n$。
\end{definition}

\begin{remark}
若 $k=\operatorname{ind}(A)$，则
\begin{equation}
 \F^n=\Image A^k\oplus\Ker A^k.
\end{equation}
事实上两者维数之和为 $n$，而若 $x=A^ky$ 同时满足 $A^kx=0$，则 $A^{2k}y=0$。稳定性给出 $\Ker A^{2k}=\Ker A^k$，故 $x=A^ky=0$。
\end{remark}

\subsection{单侧逆与满秩}

\begin{theorem}
设 $A\in\F^{m\times n}$。
\begin{enumerate}
  \item $A$ 满行秩，即 $\rank A=m$，当且仅当存在 $D\in\F^{n\times m}$ 使
  \begin{equation}
   AD=I_m.
  \end{equation}
  此时 $D$ 称为 $A$ 的右逆。
  \item $A$ 满列秩，即 $\rank A=n$，当且仅当存在 $C\in\F^{n\times m}$ 使
  \begin{equation}
   CA=I_n.
  \end{equation}
  此时 $C$ 称为 $A$ 的左逆。
\end{enumerate}
\end{theorem}

\begin{proof}
若 $A$ 满行秩，则线性映射 $A:\F^n\to\F^m$ 满射。对标准基 $e_1,\ldots,e_m$ 分别选取 $d_j$ 满足 $Ad_j=e_j$，令
\begin{equation}
 D=(d_1,\ldots,d_m),
\end{equation}
便有 $AD=I_m$。反之，$AD=I_m$ 给出
\begin{equation}
 m=\rank I_m\leq\rank A\leq m.
\end{equation}
左逆的情形可对转置使用同一论证，或把单射映射延拓为一组基后直接构造。
\end{proof}

\begin{center}
\begin{tabular}{p{0.43\textwidth}p{0.43\textwidth}}
\toprule
满行秩 $\rank A=m$ & 满列秩 $\rank A=n$\\
\midrule
行向量线性无关 & 列向量线性无关\\
$\Image A=\F^m$ & $\Ker A=\{0\}$\\
对每个 $b$，$Ax=b$ 有解 & 对每个 $b$，$Ax=b$ 至多一解\\
存在右逆 $D$，使 $AD=I_m$ & 存在左逆 $C$，使 $CA=I_n$\\
\bottomrule
\end{tabular}
\end{center}

\begin{remark}
当 $m=n$ 时，左逆、右逆、满行秩、满列秩和可逆性全部等价，而且左逆与右逆必相同。矩形矩阵则可能只有一种单侧逆。
\end{remark}

\section{习题课 \#7（10月25日）}

\subsection{Sylvester 秩不等式及等号条件}

设
\begin{equation}
 A\in\F^{m\times n},
 \qquad
 B\in\F^{n\times p}.
\end{equation}

\begin{theorem}[Sylvester 秩不等式]
有
\begin{equation}
 \rank(AB)\geq\rank A+\rank B-n.
\end{equation}
并且等号成立当且仅当
\begin{equation}
 \Ker A\subseteq\Image B.
\end{equation}
\end{theorem}

\begin{proof}
将 $A$ 限制到 $\Image B$ 上。其像为 $\Image(AB)$，核为 $\Ker A\cap\Image B$，故
\begin{equation}
 \rank(AB)
 =\rank B-\dim(\Ker A\cap\Image B).
\end{equation}
因为
\begin{equation}
 \dim(\Ker A\cap\Image B)\leq\dim\Ker A=n-\rank A,
\end{equation}
即得不等式。等号成立恰好要求
\begin{equation}
 \dim(\Ker A\cap\Image B)=\dim\Ker A,
\end{equation}
也就是 $\Ker A\subseteq\Image B$。
\end{proof}

\begin{corollary}
下列条件等价：
\begin{enumerate}
  \item $\rank(AB)=\rank A+\rank B-n$；
  \item $\Ker A\subseteq\Image B$；
  \item 存在 $X\in\F^{n\times m}$ 与 $Y\in\F^{p\times n}$，使
  \begin{equation}
   XA+BY=I_n.
  \end{equation}
\end{enumerate}
\end{corollary}

\begin{proof}
第三个条件意味着每个 $v\in\Ker A$ 满足 $v=BYv\in\Image B$。反之，若 $\Ker A\subseteq\Image B$，取直和分解
\begin{equation}
 \F^n=U\oplus\Ker A.
\end{equation}
映射 $A|_U$ 单射。可构造 $X$ 使 $XA$ 是到 $U$ 的投影，再构造 $Y$ 使 $BY$ 是到 $\Ker A$ 的投影，于是二者之和为恒等映射。
\end{proof}

\subsection{满秩分解与秩标准形}

\begin{theorem}[满秩分解]
若 $A\in\F^{m\times n}$ 且 $\rank A=r$，则存在
\begin{equation}
 P\in\F^{m\times r},
 \qquad
 Q\in\F^{r\times n}
\end{equation}
使
\begin{equation}
 A=PQ,
 \qquad
 \rank P=\rank Q=r.
\end{equation}
\end{theorem}

\begin{proof}
从 $A$ 中选取 $r$ 个极大线性无关列组成 $P$。$A$ 的每一列都可唯一表示为 $P$ 的列向量的线性组合，把这些坐标作为 $Q$ 的列便有 $A=PQ$。由于 $A$ 的秩为 $r$，$Q$ 必满行秩。
\end{proof}

\begin{proposition}[满秩分解的唯一性]
若
\begin{equation}
 A=P_1Q_1=P_2Q_2
\end{equation}
都是满秩分解，则存在唯一 $S\in\GL_r(\F)$ 使
\begin{equation}
 P_2=P_1S,
 \qquad
 Q_2=S^{-1}Q_1.
\end{equation}
\end{proposition}

\begin{proof}
$P_1,P_2$ 的列空间都等于 $\Image A$，故 $P_2=P_1S$。两组列都是基，所以 $S$ 可逆。代回并利用 $P_1$ 满列秩即可得到 $Q_2=S^{-1}Q_1$。
\end{proof}

\begin{corollary}[秩标准形]
存在可逆矩阵 $U\in\GL_m(\F)$、$V\in\GL_n(\F)$，使
\begin{equation}
 UAV=\begin{pmatrix}
 I_r&0\\
 0&0
 \end{pmatrix}.
\end{equation}
因此两个同型矩阵等价当且仅当它们的秩相同。
\end{corollary}

\begin{remark}
若经过行列置换后，$A$ 的左上角 $r\times r$ 子矩阵 $A_{11}$ 可逆，写成
\begin{equation}
 A=\begin{pmatrix}
 A_{11}&A_{12}\\
 A_{21}&A_{22}
 \end{pmatrix},
\end{equation}
则 $\rank A=r$ 当且仅当
\begin{equation}
 A_{22}=A_{21}A_{11}^{-1}A_{12}.
\end{equation}
这就是 Schur 补为零的情形。
\end{remark}

\subsection{矩阵和的秩}

\begin{theorem}
对同型矩阵 $A,B$，
\begin{equation}
 \rank(A+B)\leq\rank A+\rank B.
\end{equation}
等号成立当且仅当
\begin{equation}
 \Image A\cap\Image B=\{0\}
\end{equation}
且
\begin{equation}
 \Image A^T\cap\Image B^T=\{0\}.
\end{equation}
复数情形可将转置理解为行空间，而不必使用共轭转置。
\end{theorem}

\begin{proof}
因为
\begin{equation}
 \Image(A+B)\subseteq\Image A+\Image B,
\end{equation}
先得上界。若等号成立，则
\begin{equation}
 \rank(A+B)=\dim(\Image A+\Image B)
\end{equation}
迫使列空间之交为零；对转置应用同一论证得到行空间之交为零。

反之，若两种交都为零，可选取适当的可逆行、列变换，使
\begin{equation}
 A\sim
 \begin{pmatrix}
 I_r&0&0\\
 0&0&0\\
 0&0&0
 \end{pmatrix},
 \qquad
 B\sim
 \begin{pmatrix}
 0&0&0\\
 0&I_s&0\\
 0&0&0
 \end{pmatrix},
\end{equation}
且两者使用相同的外侧变换。于是 $A+B$ 的秩为 $r+s$。
\end{proof}

\begin{corollary}
对同型矩阵 $A,B$，以下三条中的任意两条推出第三条：
\begin{equation}
 \rank(A+B)=\rank A+\rank B,
\end{equation}
\begin{equation}
 \rank(A^T,B^T)=\rank A+\rank B,
\end{equation}
以及相应列空间、行空间直和条件。
\end{corollary}

\subsection{Binet--Cauchy 公式与伴随矩阵}

\begin{theorem}[Binet--Cauchy 公式]
设 $A\in\F^{m\times n}$、$B\in\F^{n\times m}$，且 $m\leq n$。对 $m$ 元指标集
\begin{equation}
 I=\{i_1<\cdots<i_m\}\subseteq\{1,\ldots,n\},
\end{equation}
记 $A_{:,I}$ 为取 $A$ 的这些列所得方阵，$B_{I,:}$ 为取 $B$ 的这些行所得方阵。则
\begin{equation}
 \det(AB)=\sum_{|I|=m}\det(A_{:,I})\det(B_{I,:}).
\end{equation}
\end{theorem}

\begin{proof}
把 $AB$ 的每一列写成 $A$ 的列向量的线性组合，并对行列式使用多重线性性。若所选指标有重复，相应行列式为零；将互异指标按递增顺序重排后，两个置换符号相消，便得到上述求和式。
\end{proof}

\begin{corollary}[Gram 行列式]
若 $A\in\R^{n\times m}$ 且 $n\geq m$，则
\begin{equation}
 \det(A^TA)=\sum_{|I|=m}\det(A_{I,:})^2\geq0.
\end{equation}
并且 $\det(A^TA)>0$ 当且仅当 $A$ 满列秩。
\end{corollary}

\begin{proposition}
对任意同阶方阵 $A,B$，有
\begin{equation}
 \adj(AB)=\adj(B)\adj(A).
\end{equation}
\end{proposition}

\begin{proof}
当 $A,B$ 都可逆时，
\begin{equation}
 \adj(AB)=\det(AB)(AB)^{-1}
 =\det B\,B^{-1}\det A\,A^{-1}
 =\adj(B)\adj(A).
\end{equation}
两边的每个矩阵元都是 $A,B$ 各元素的多项式。把 $A,B$ 分别替换为 $A+tI,B+sI$，在除有限多个参数值外二者可逆，于是多项式恒等式对所有 $A,B$ 成立。
\end{proof}

\subsection{低秩修正与 Schur 补}

\begin{theorem}[对角矩阵加秩一矩阵]
设
\begin{equation}
 D=\diag(a_1,\ldots,a_n)
\end{equation}
可逆，且
\begin{equation}
 1+\ones^TD^{-1}\ones\neq0.
\end{equation}
则
\begin{equation}
 (D+\ones\ones^T)^{-1}
 =D^{-1}-
 \frac{D^{-1}\ones\ones^TD^{-1}}
 {1+\ones^TD^{-1}\ones}.
\end{equation}
\end{theorem}

\begin{proof}
将右端记为 $R$，直接计算
\begin{equation}
 (D+\ones\ones^T)R=I_n.
\end{equation}
这也是 Sherman--Morrison 公式的特例。
\end{proof}

\begin{theorem}[Schur 补与分块逆]
设
\begin{equation}
 M=\begin{pmatrix}A&B\\ C&D\end{pmatrix},
\end{equation}
其中 $A$ 可逆，并记
\begin{equation}
 S=D-CA^{-1}B.
\end{equation}
则
\begin{equation}
 \begin{pmatrix}I&0\\-CA^{-1}&I\end{pmatrix}
 M
 \begin{pmatrix}I&-A^{-1}B\\0&I\end{pmatrix}
 =\begin{pmatrix}A&0\\0&S\end{pmatrix}.
\end{equation}
因此 $M$ 可逆当且仅当 $S$ 可逆，并且
\begin{equation}
 M^{-1}=
 \begin{pmatrix}
 A^{-1}+A^{-1}BS^{-1}CA^{-1}&-A^{-1}BS^{-1}\\
 -S^{-1}CA^{-1}&S^{-1}
 \end{pmatrix}.
\end{equation}
\end{theorem}

\begin{corollary}
若 $M=M^T$ 且 $A$ 正定，则
\begin{equation}
 M\text{ 正定}
 \quad\Longleftrightarrow\quad
 S=D-B^TA^{-1}B\text{ 正定}.
\end{equation}
\end{corollary}

\subsection{Vandermonde 矩阵的逆}

令 $x_1,\ldots,x_n$ 两两不同，并取标准 Vandermonde 矩阵
\begin{equation}
 V=(x_i^{j-1})_{i,j=1}^n.
\end{equation}
定义 Lagrange 基函数
\begin{equation}
 \ell_i(t)=\prod_{k\neq i}\frac{t-x_k}{x_i-x_k}
 =\sum_{j=1}^n c_{ji}t^{j-1}.
\end{equation}
由于 $\ell_i(x_k)=\delta_{ik}$，有
\begin{equation}
 V(c_{ji})_{j,i=1}^n=I_n.
\end{equation}

\begin{theorem}
Vandermonde 矩阵的逆满足
\begin{equation}
 (V^{-1})_{ji}=[t^{j-1}]\ell_i(t),
\end{equation}
即第 $i$ 列由第 $i$ 个 Lagrange 基多项式的系数组成。显式地，
\begin{equation}
 (V^{-1})_{ji}
 =\frac{(-1)^{n-j}
 e_{n-j}(x_1,\ldots,\widehat{x_i},\ldots,x_n)}
 {\prod_{k\neq i}(x_i-x_k)},
\end{equation}
其中 $e_k$ 是初等对称多项式，帽号表示删去该变量。
\end{theorem}

\begin{remark}[单位根情形]
取
\begin{equation}
 \omega=\e^{-2\pi\ii/n},
 \qquad
 F=(\omega^{jk})_{j,k=0}^{n-1},
\end{equation}
则几何级数给出
\begin{equation}
 F^*F=nI_n,
 \qquad
 F^{-1}=\frac1nF^*.
\end{equation}
这正是 Vandermonde 逆在单位根节点上的特殊简化。
\end{remark}

\subsection{行列式交换恒等式与矩阵逆公式}

\begin{theorem}[Sylvester 行列式恒等式]
若 $A\in\F^{m\times n}$、$B\in\F^{n\times m}$，则
\begin{equation}
 \det(I_m+AB)=\det(I_n+BA).
\end{equation}
更一般地，作为关于 $\lambda$ 的多项式，
\begin{equation}
 \det(\lambda I_m-AB)
 =\lambda^{m-n}\det(\lambda I_n-BA)
\end{equation}
在 $m\geq n$ 时成立；若 $m<n$，交换 $m,n$ 后理解该式。
\end{theorem}

\begin{proof}
对分块矩阵
\begin{equation}
 \begin{pmatrix}I_m&A\\-B&I_n\end{pmatrix}
\end{equation}
分别以左上角和右下角作 Schur 补，即得第一式。第二式先在 $\lambda\neq0$ 时提出 $\lambda$ 并应用第一式，再由多项式恒等性延拓到所有 $\lambda$。
\end{proof}

\begin{corollary}
$AB$ 与 $BA$ 的非零特征值连同代数重数完全相同。
\end{corollary}

\begin{proposition}[Neumann 级数]
若某个次乘法矩阵范数满足 $\norm{A}<1$，则
\begin{equation}
 (I-A)^{-1}=I+A+A^2+\cdots,
\end{equation}
级数在该范数下绝对收敛。更一般地，只要谱半径 $\rho(A)<1$，该结论仍成立。
\end{proposition}

\begin{proposition}[推入恒等式]
若相关逆存在，则
\begin{equation}
 (I-AB)^{-1}=I+A(I-BA)^{-1}B.
\end{equation}
等价地，
\begin{equation}
 (I-AB)^{-1}A=A(I-BA)^{-1}.
\end{equation}
\end{proposition}

\begin{theorem}[Woodbury 公式]
若各逆均存在，则
\begin{equation}
 (A+UCV)^{-1}
 =A^{-1}-A^{-1}U
 (C^{-1}+VA^{-1}U)^{-1}
 VA^{-1}.
\end{equation}
当 $C=1$ 且 $U=u,V=v^T$ 时，得到 Sherman--Morrison 公式
\begin{equation}
 (A+uv^T)^{-1}
 =A^{-1}-\frac{A^{-1}uv^TA^{-1}}
 {1+v^TA^{-1}u}.
\end{equation}
\end{theorem}

\subsection{DFT、循环矩阵与快速算法}

取 $\omega=\e^{-2\pi\ii/n}$，离散 Fourier 矩阵为
\begin{equation}
 F=(\omega^{jk})_{j,k=0}^{n-1}.
\end{equation}
对向量 $x\in\C^n$，其离散 Fourier 变换和逆变换分别为
\begin{equation}
 \widehat{x}=Fx,
 \qquad
 x=\frac1nF^*\widehat{x}.
\end{equation}

\begin{theorem}[循环矩阵的对角化]
设 $C$ 的第一列为 $c=(c_0,\ldots,c_{n-1})^T$，其余列由循环移位得到。则
\begin{equation}
 C=F^{-1}\diag(Fc)F.
\end{equation}
因此循环卷积满足卷积定理
\begin{equation}
 F(x\ast y)=(Fx)\odot(Fy),
\end{equation}
其中 $\odot$ 表示逐分量乘积。
\end{theorem}

\begin{proof}
Fourier 向量
\begin{equation}
 v_k=(1,\omega^k,\omega^{2k},\ldots,\omega^{(n-1)k})^T
\end{equation}
是每个循环移位算子的特征向量，因而也是任意循环矩阵的特征向量。相应特征值正是 $Fc$ 的各分量。
\end{proof}

\begin{remark}
直接计算 DFT 需要 $O(n^2)$ 次运算。快速 Fourier 变换利用 $n$ 的因子分解递归拆分偶、奇指标，将复杂度降为 $O(n\log n)$。把单位根换成有限域中的适当本原根便得到数论变换，可用于避免浮点误差的多项式乘法。
\end{remark}

\section{习题课 \#8（11月1日）}

\subsection{Gram--Schmidt 正交化}

设 $V$ 是实或复内积空间，$\alpha_1,\ldots,\alpha_s$ 线性无关。Gram--Schmidt 过程递归定义
\begin{align}
 \widetilde\beta_1&=\alpha_1,
 &\beta_1&=\frac{\widetilde\beta_1}{\norm{\widetilde\beta_1}},\\
 \widetilde\beta_k&=\alpha_k-
 \sum_{j=1}^{k-1}\inner{\beta_j}{\alpha_k}\beta_j,
 &\beta_k&=\frac{\widetilde\beta_k}{\norm{\widetilde\beta_k}},
 \qquad k=2,\ldots,s.
\end{align}
复内积取第一变量共轭线性、第二变量线性时，应按所采用的约定相应调整内积次序。

\begin{theorem}
上述过程满足
\begin{equation}
 \inner{\beta_i}{\beta_j}=\delta_{ij},
 \qquad
 \Span(\beta_1,\ldots,\beta_k)
 =\Span(\alpha_1,\ldots,\alpha_k)
\end{equation}
对每个 $k$ 成立。因此 $\beta_1,\ldots,\beta_s$ 是与原向量组张成同一子空间的一组标准正交基。
\end{theorem}

\begin{proof}
$\widetilde\beta_k$ 是从 $\alpha_k$ 中减去其在此前标准正交向量上的投影，故与 $\beta_1,\ldots,\beta_{k-1}$ 正交。若 $\widetilde\beta_k=0$，则 $\alpha_k$ 属于前 $k-1$ 个向量的张成空间，与线性无关矛盾。归一化后即得结论。
\end{proof}

\begin{remark}[由坐标变换诱导的内积]
若 $A=(\alpha_1,\ldots,\alpha_n)$ 可逆，向量 $x$ 表示在基 $\alpha_1,\ldots,\alpha_n$ 下的坐标，则
\begin{equation}
 \norm{Ax}_2^2=x^TA^TAx,
 \qquad
 \inner{Ax}{Ay}=x^TA^TAy.
\end{equation}
因此 $A^TA$ 是这组基对应的 Gram 矩阵。距离与夹角都可在坐标空间中通过该正定矩阵计算。
\end{remark}

\subsection{QR 分解}

\begin{theorem}[薄 QR 分解]
设 $A\in\R^{m\times n}$，$m\geq n$，且 $A$ 满列秩。则存在
\begin{equation}
 Q\in\R^{m\times n},
 \qquad
 R\in\R^{n\times n}
\end{equation}
使
\begin{equation}
 A=QR,
 \qquad Q^TQ=I_n,
\end{equation}
并且 $R$ 为对角元全正的上三角矩阵。这样的分解唯一。
\end{theorem}

\begin{proof}
对 $A$ 的列向量作 Gram--Schmidt 正交化，把所得标准正交向量作为 $Q$ 的列。每个原列都由当前及此前的 $Q$ 列线性表示，故系数矩阵 $R$ 为上三角；其对角元是未归一化残差的范数，因而为正。

若 $A=Q_1R_1=Q_2R_2$ 都满足条件，则
\begin{equation}
 Q_2^TQ_1=R_2R_1^{-1}.
\end{equation}
左端是正交矩阵，右端是对角元全正的上三角矩阵。唯一可能是恒等矩阵，故 $Q_1=Q_2$、$R_1=R_2$。
\end{proof}

\begin{remark}
若 $A$ 只有秩 $r<n$，仍可对一组极大无关列作薄 QR 分解，但不再能要求一个 $n\times n$ 的 $R$ 对角元全非零。数值计算中常配合列主元选取写成
\begin{equation}
 A\Pi=QR.
\end{equation}
\end{remark}

\subsection{最小二乘问题}

考虑
\begin{equation}
 \min_{x\in\R^n}\norm{Ax-b}_2,
 \qquad A\in\R^{m\times n}.
\end{equation}

\begin{theorem}[正规方程]
向量 $\widehat x$ 是最小二乘解当且仅当
\begin{equation}
 A^T(A\widehat x-b)=0,
\end{equation}
即
\begin{equation}
 A^TA\widehat x=A^Tb.
\end{equation}
若 $A$ 满列秩，则解唯一，并且
\begin{equation}
 \widehat x=(A^TA)^{-1}A^Tb.
\end{equation}
\end{theorem}

\begin{proof}
最小值处残差 $r=b-A\widehat x$ 必与 $\Image A$ 正交，故 $A^Tr=0$。反之，若 $A^Tr=0$，则对任意 $x$，
\begin{equation}
 \norm{Ax-b}_2^2
 =\norm{A(x-\widehat x)-r}_2^2
 =\norm{A(x-\widehat x)}_2^2+\norm r_2^2,
\end{equation}
所以 $\widehat x$ 确为最小点。
\end{proof}

\begin{corollary}[QR 求解]
若 $A=QR$ 是薄 QR 分解，则
\begin{equation}
 \widehat x=R^{-1}Q^Tb,
\end{equation}
并且最小残差平方为
\begin{equation}
 \norm b_2^2-\norm{Q^Tb}_2^2.
\end{equation}
正交投影到 $\Image A$ 的矩阵为
\begin{equation}
 P=QQ^T=A(A^TA)^{-1}A^T.
\end{equation}
\end{corollary}

\subsection{秩不等式汇总}

对尺寸相容的矩阵，有以下常用估计：
\begin{align}
 \rank(A+B)&\leq\rank A+\rank B,\\
 \rank(A,B)&\leq\rank A+\rank B,\\
 \rank\begin{pmatrix}A\\B\end{pmatrix}
 &\leq\rank A+\rank B,\\
 \rank(AB)&\leq\min\{\rank A,\rank B\},\\
 \rank(AB)&\geq\rank A+\rank B-n,\\
 \rank(ABC)&\geq\rank(AB)+\rank(BC)-\rank B.
\end{align}
最后一条称为 Frobenius 秩不等式。

\begin{proof}[Frobenius 不等式的证明]
取 $B=P\begin{psmallmatrix}I_r&0\\0&0\end{psmallmatrix}Q$ 的秩标准形，把外侧可逆矩阵吸收到 $A,C$ 中。此时 $AB$ 只依赖 $A$ 的前 $r$ 列，$BC$ 只依赖 $C$ 的前 $r$ 行，而 $ABC$ 是这两个截取矩阵的乘积。对它们应用 Sylvester 秩不等式即得
\begin{equation}
 \rank(ABC)\geq\rank(AB)+\rank(BC)-r.
\end{equation}
这里 $r=\rank B$。
\end{proof}

\begin{exercise}
设同阶方阵满足
\begin{equation}
 A+B=AB.
\end{equation}
证明 $A$ 与 $B$ 可交换，并写出它们与 $I$ 的关系。
\end{exercise}

\begin{solution}
原式等价于
\begin{equation}
 (I-A)(I-B)=I.
\end{equation}
方阵的右逆也是左逆，故
\begin{equation}
 (I-B)(I-A)=I.
\end{equation}
展开得
\begin{equation}
 A+B=BA.
\end{equation}
与已知式比较可知 $AB=BA$，并且 $I-A$ 与 $I-B$ 互为逆矩阵。
\end{solution}

\subsection{Gram 行列式与分块估计}

设 $A=(B,C)$，其中 $B$、$C$ 分别由一部分列组成。Gram 矩阵为
\begin{equation}
 A^TA=\begin{pmatrix}
 B^TB&B^TC\\
 C^TB&C^TC
\end{pmatrix}.
\end{equation}

\begin{theorem}[Gram--Hadamard/Fischer 型不等式]
有
\begin{equation}
 \det(A^TA)
 \leq\det(B^TB)\det(C^TC).
\end{equation}
若 $B,C$ 均满列秩，则等号成立当且仅当
\begin{equation}
 B^TC=0,
\end{equation}
即两个列空间正交。
\end{theorem}

\begin{proof}
若 $B$ 满列秩，对左上角作 Schur 补：
\begin{equation}
 \det(A^TA)
 =\det(B^TB)
 \det\bigl(C^TC-C^TB(B^TB)^{-1}B^TC\bigr).
\end{equation}
括号中等于
\begin{equation}
 C^T(I-P_B)C,
 \qquad
 P_B=B(B^TB)^{-1}B^T,
\end{equation}
且满足
\begin{equation}
 0\preccurlyeq C^T(I-P_B)C\preccurlyeq C^TC.
\end{equation}
比较特征值乘积即得不等式。退化情形可由连续性得到。
\end{proof}

\subsection{内积空间中的基本不等式}

\begin{theorem}[勾股定理]
若 $\alpha\perp\beta$，则
\begin{equation}
 \norm{\alpha+\beta}^2=\norm\alpha^2+\norm\beta^2.
\end{equation}
\end{theorem}

\begin{theorem}[Cauchy--Schwarz 不等式]
对任意 $\alpha,\beta$，
\begin{equation}
 \abs{\inner\alpha\beta}\leq\norm\alpha\,\norm\beta.
\end{equation}
等号成立当且仅当 $\alpha,\beta$ 线性相关。
\end{theorem}

\begin{proof}
若 $\beta\neq0$，考虑
\begin{equation}
 0\leq\norm{\alpha-t\beta}^2.
\end{equation}
在实数情形取
\begin{equation}
 t=\frac{\inner\beta\alpha}{\norm\beta^2},
\end{equation}
在复数情形取相应的正交投影系数，即得
\begin{equation}
 0\leq\norm\alpha^2-
 \frac{\abs{\inner\alpha\beta}^2}{\norm\beta^2}.
\end{equation}
\end{proof}

\begin{corollary}[三角不等式]
\begin{equation}
 \norm{\alpha+\beta}\leq\norm\alpha+\norm\beta.
\end{equation}
\end{corollary}

\subsection{平行四边形法则与极化恒等式}

\begin{theorem}[平行四边形法则]
内积诱导的范数满足
\begin{equation}
 \norm{\alpha+\beta}^2+\norm{\alpha-\beta}^2
 =2\bigl(\norm\alpha^2+\norm\beta^2\bigr).
\end{equation}
\end{theorem}

\begin{theorem}[极化恒等式]
在实内积空间中，
\begin{equation}
 \inner\alpha\beta
 =\frac14\left(
 \norm{\alpha+\beta}^2-\norm{\alpha-\beta}^2
 \right).
\end{equation}
在复内积空间中，若内积对第二个变量线性，则
\begin{equation}
 \inner\alpha\beta
 =\frac14\left(
 \norm{\alpha+\beta}^2-\norm{\alpha-\beta}^2
 -\ii\norm{\alpha+\ii\beta}^2
 +\ii\norm{\alpha-\ii\beta}^2
 \right).
\end{equation}
\end{theorem}

\begin{remark}
平行四边形法则不仅是内积范数的必要条件，也是充分条件。一个范数满足该法则时，可由极化公式恢复唯一的内积。
\end{remark}

\section{习题课 \#9（11月8日）}

\subsection{线性映射的核、像与矩阵表示}

设 $V_1,V_2$ 是域 $\F$ 上的向量空间。映射 $T:V_1\to V_2$ 称为线性映射，如果
\begin{equation}
 T(au+bv)=aT(u)+bT(v),
 \qquad a,b\in\F,
 \quad u,v\in V_1.
\end{equation}
其核与像分别为
\begin{equation}
 \Ker T=\set{v\in V_1:T(v)=0},
 \qquad
 \Image T=\set{T(v):v\in V_1}.
\end{equation}
$T$ 单射当且仅当 $\Ker T=\{0\}$，满射当且仅当 $\Image T=V_2$。

固定 $V_1,V_2$ 的基后，线性映射由一个矩阵表示。若更换定义域与值域的基，表示矩阵在左右分别乘以可逆矩阵，因此同一线性映射的不同矩阵表示彼此等价。

\begin{example}
令
\begin{equation}
 A=\begin{pmatrix}
 3&1&4&7\\
 -1&1&0&3
 \end{pmatrix}.
\end{equation}
前两列线性无关，后两列可由它们表示，所以
\begin{equation}
 \Image A
 =\Span\left\{
 \begin{pmatrix}3\\-1\end{pmatrix},
 \begin{pmatrix}1\\1\end{pmatrix}
 \right\},
 \qquad
 \rank A=2.
\end{equation}
由 $\dim\Ker A=4-2=2$，消元可得
\begin{equation}
 \Ker A=\Span\left\{
 \begin{pmatrix}-1\\-1\\1\\0\end{pmatrix},
 \begin{pmatrix}-1\\-4\\0\\1\end{pmatrix}
 \right\}.
\end{equation}
\end{example}

\subsection{矩阵方程 AXB=C}

先考虑可逆情形。若 $A,B$ 都可逆，则
\begin{equation}
 AXB=C
 \quad\Longleftrightarrow\quad
 X=A^{-1}CB^{-1}.
\end{equation}
一般情形可使用广义逆统一描述。

\begin{definition}
矩阵 $A^{-}$ 称为 $A$ 的一个内逆，如果
\begin{equation}
 AA^{-}A=A.
\end{equation}
每个有限维矩阵都存在内逆。事实上，若
\begin{equation}
 A=P\begin{pmatrix}I_r&0\\0&0\end{pmatrix}Q
\end{equation}
是秩标准形，则可取
\begin{equation}
 A^{-}=Q^{-1}
 \begin{pmatrix}I_r&0\\0&0\end{pmatrix}
 P^{-1},
\end{equation}
其中零块按尺寸补齐。
\end{definition}

\begin{theorem}[矩阵方程的相容条件与通解]
设
\begin{equation}
 A\in\F^{m\times n},
 \quad X\in\F^{n\times p},
 \quad B\in\F^{p\times q},
 \quad C\in\F^{m\times q},
\end{equation}
并任选内逆 $A^{-},B^{-}$。方程
\begin{equation}
 AXB=C
\end{equation}
有解当且仅当
\begin{equation}
 AA^{-}CB^{-}B=C.
\end{equation}
相容时，全部解为
\begin{equation}
 X=A^{-}CB^{-}+Y-A^{-}AYBB^{-},
 \qquad Y\in\F^{n\times p}\text{ 任意}.
\end{equation}
\end{theorem}

\begin{proof}
若 $AXB=C$，则
\begin{equation}
 AA^{-}CB^{-}B
 =AA^{-}AXBB^{-}B
 =AXB=C.
\end{equation}
反之，相容条件成立时，取 $X_0=A^{-}CB^{-}$ 即有 $AX_0B=C$。

对任意 $Y$，代入通解表达式得
\begin{align}
 AXB
 &=AA^{-}CB^{-}B+AYB-AA^{-}AYBB^{-}B\\
 &=C+AYB-AYB=C.
\end{align}
若 $X$ 是任一解，在右端取 $Y=X$，则
\begin{equation}
 A^{-}CB^{-}+X-A^{-}AXBB^{-}=X,
\end{equation}
故该公式确实包含全部解。
\end{proof}

\begin{remark}
相容条件可解释为两个空间条件：$C$ 的每一列必须属于 $\Image A$，而 $C$ 的每一行必须属于 $B$ 的行空间。若只求 $AX=C$，则对 $C$ 的每一列分别解同一个系数矩阵的方程即可；若只求 $XB=C$，可转置后同样处理。
\end{remark}

\subsection{秩标准形中的像与核}

设 $\rank A=r$，并写成
\begin{equation}
 A=P
 \begin{pmatrix}I_r&0\\0&0\end{pmatrix}
 Q,
 \qquad P,Q\text{ 可逆}.
\end{equation}
把
\begin{equation}
 P=(P_1,P_2),
 \qquad
 Q^{-1}=(Q_1,Q_2)
\end{equation}
按秩分块，则
\begin{equation}
 \Image A=\Image P_1,
 \qquad
 \Ker A=\Image Q_2.
\end{equation}
相应地，
\begin{equation}
 \Image A^T=\Image Q^T_1,
 \qquad
 \Ker A^T=\Image P^{-T}_2.
\end{equation}
这些公式把像空间、零空间和秩标准形直接联系起来。

\begin{proposition}
若 $A=P\begin{psmallmatrix}I_r&0\\0&0\end{psmallmatrix}Q$，则
\begin{equation}
 \Ker A=\Ker\begin{pmatrix}I_r&0\\0&0\end{pmatrix}Q,
\end{equation}
而左乘可逆矩阵不改变核，右乘可逆矩阵只作坐标变换。类似地，右乘可逆矩阵不改变列空间的维数，左乘可逆矩阵把列空间同构地变换。
\end{proposition}

\subsection{分块矩阵与 Schur 补}

\begin{theorem}
设
\begin{equation}
 M=\begin{pmatrix}A_{11}&A_{12}\\A_{21}&A_{22}\end{pmatrix}
\end{equation}
且 $A_{11}$ 可逆。记
\begin{equation}
 S=A_{22}-A_{21}A_{11}^{-1}A_{12}.
\end{equation}
则 $M$ 可逆当且仅当 $S$ 可逆，并且
\begin{equation}
 M^{-1}=\begin{pmatrix}
 A_{11}^{-1}+A_{11}^{-1}A_{12}S^{-1}A_{21}A_{11}^{-1}
 &-A_{11}^{-1}A_{12}S^{-1}\\
 -S^{-1}A_{21}A_{11}^{-1}&S^{-1}
 \end{pmatrix}.
\end{equation}
\end{theorem}

\begin{proof}
利用分解
\begin{equation}
 M=
 \begin{pmatrix}I&0\\A_{21}A_{11}^{-1}&I\end{pmatrix}
 \begin{pmatrix}A_{11}&0\\0&S\end{pmatrix}
 \begin{pmatrix}I&A_{11}^{-1}A_{12}\\0&I\end{pmatrix}
\end{equation}
即可。
\end{proof}

\subsection{若干秩恒等式与可解性判别}

\begin{theorem}
设 $A\in\F^{m\times n}$、$B\in\F^{n\times m}$。则
\begin{equation}
 \rank(A-ABA)=\rank A+\rank(I_n-BA)-n.
\end{equation}
\end{theorem}

\begin{proof}
写成
\begin{equation}
 A-ABA=A(I_n-BA).
\end{equation}
Sylvester 秩不等式给出下界。另一方面，若 $x\in\Ker A$，则
\begin{equation}
 (I_n-BA)x=x,
\end{equation}
所以 $\Ker A\subseteq\Image(I_n-BA)$。这正是 Sylvester 不等式取等号的条件。
\end{proof}

\begin{proposition}
设矩阵尺寸相容。方程
\begin{equation}
 ABX=A
\end{equation}
有解当且仅当
\begin{equation}
 \rank(AB)=\rank A.
\end{equation}
\end{proposition}

\begin{proof}
若有解，则 $\Image A\subseteq\Image(AB)$，而反向包含总成立，故两列空间相同。反之，若秩相同，则 $\Image A=\Image(AB)$，逐列求解即可构造 $X$。
\end{proof}

\begin{proposition}
对同阶方阵 $A,B$，
\begin{equation}
 \rank(I-AB)\leq\rank(I-A)+\rank(I-B).
\end{equation}
\end{proposition}

\begin{proof}
由
\begin{equation}
 I-AB=(I-A)+A(I-B)
\end{equation}
及矩阵和、矩阵乘积的秩上界立即得到。
\end{proof}

\begin{theorem}
若同阶方阵满足
\begin{equation}
 AB=BA=0,
 \qquad
 \rank A^2=\rank A,
\end{equation}
则
\begin{equation}
 \rank(A+B)=\rank A+\rank B.
\end{equation}
\end{theorem}

\begin{proof}
条件 $\rank A^2=\rank A$ 等价于
\begin{equation}
 \Image A\cap\Ker A=\{0\}.
\end{equation}
由 $AB=0$ 得 $\Image B\subseteq\Ker A$，所以
\begin{equation}
 \Image A\cap\Image B=\{0\}.
\end{equation}
对转置矩阵使用 $BA=0$ 与 $\rank(A^T)^2=\rank A^T$，同样得到行空间之交为零。由矩阵和秩等号的判别即得结论。
\end{proof}

\subsection{LU 分解}

\begin{theorem}
设 $A\in\F^{n\times n}$ 的各阶顺序主子式均非零：
\begin{equation}
 \det A_{1:k,1:k}\neq0,
 \qquad k=1,\ldots,n.
\end{equation}
则存在唯一分解
\begin{equation}
 A=LU,
\end{equation}
其中 $L$ 是对角元全为 $1$ 的下三角矩阵，$U$ 是可逆上三角矩阵。
\end{theorem}

\begin{proof}
顺序主子式非零保证 Gaussian 消元无需换行，所有消元乘子组成 $L$，消元后的阶梯矩阵为 $U$。若另有 $A=\widetilde L\widetilde U$，则
\begin{equation}
 \widetilde L^{-1}L=\widetilde U U^{-1}.
\end{equation}
左端是单位下三角，右端是上三角，只能都等于 $I$。
\end{proof}

\begin{remark}
若允许置换，一般可写成
\begin{equation}
 PA=LU.
\end{equation}
数值计算中选主元的目的不仅是避免零主元，也用于控制舍入误差。
\end{remark}

\subsection{Householder 变换与 QR 分解}

\begin{definition}
对单位向量 $w\in\R^n$，矩阵
\begin{equation}
 H=I-2ww^T
\end{equation}
称为 Householder 反射。
\end{definition}

\begin{proposition}
Householder 矩阵满足
\begin{equation}
 H^T=H,
 \qquad
 H^2=I,
 \qquad
 H^TH=I.
\end{equation}
它把 $w$ 方向取反，并保持 $w^\perp$ 中的向量不变。
\end{proposition}

\begin{theorem}
给定非零向量 $x\in\R^n$，可选单位向量 $w$ 使相应 Householder 反射满足
\begin{equation}
 Hx=\pm\norm{x}_2e_1.
\end{equation}
例如取
\begin{equation}
 v=x+\operatorname{sgn}(x_1)\norm{x}_2e_1,
 \qquad
 w=\frac{v}{\norm v_2},
\end{equation}
即可避免严重的相消误差。
\end{theorem}

\begin{remark}
依次对矩阵第一列、第二列以下部分等应用 Householder 反射，可将 $A$ 化为上三角或上梯形矩阵：
\begin{equation}
 H_k\cdots H_1A=R.
\end{equation}
于是
\begin{equation}
 A=QR,
 \qquad
 Q=H_1\cdots H_k
\end{equation}
是 QR 分解。该方法比经典 Gram--Schmidt 通常更稳定。
\end{remark}

\section{习题课 \#10（11月15日）}

\subsection{矩阵的等价与线性映射的表示}

设 $T:V\to W$ 是线性映射。取 $V$ 的基
\begin{equation}
 \mathcal A=(\alpha_1,\ldots,\alpha_n),
 \qquad
 \widetilde{\mathcal A}=(\widetilde\alpha_1,\ldots,\widetilde\alpha_n),
\end{equation}
以及 $W$ 的基
\begin{equation}
 \mathcal B=(\beta_1,\ldots,\beta_m),
 \qquad
 \widetilde{\mathcal B}=(\widetilde\beta_1,\ldots,\widetilde\beta_m).
\end{equation}
若
\begin{equation}
 [v]_{\mathcal A}=Q[v]_{\widetilde{\mathcal A}},
 \qquad
 [w]_{\mathcal B}=P[w]_{\widetilde{\mathcal B}},
\end{equation}
且 $T$ 在两组基下的矩阵分别为 $A$、$\widetilde A$，则
\begin{equation}
 \widetilde A=P^{-1}AQ.
\end{equation}
因此，同一线性映射在不同定义域和值域基下的矩阵彼此等价。

\begin{definition}
同型矩阵 $A,B\in\F^{m\times n}$ 称为等价，如果存在可逆矩阵
\begin{equation}
 P\in\GL_m(\F),
 \qquad Q\in\GL_n(\F)
\end{equation}
使
\begin{equation}
 B=PAQ.
\end{equation}
\end{definition}

\begin{theorem}
两个同型矩阵等价当且仅当它们的秩相同。每个秩为 $r$ 的 $m\times n$ 矩阵都等价于
\begin{equation}
 \begin{pmatrix}I_r&0\\0&0\end{pmatrix}.
\end{equation}
\end{theorem}

\subsection{相似与线性变换}

当 $V=W$ 且定义域和值域同时更换为同一组基时，左右变换互为逆矩阵。

\begin{definition}
方阵 $A,B\in\F^{n\times n}$ 称为相似，如果存在 $Q\in\GL_n(\F)$ 使
\begin{equation}
 B=Q^{-1}AQ.
\end{equation}
相似矩阵表示同一个线性变换在不同基下的矩阵。
\end{definition}

\begin{proposition}
对任意 $Q\in\GL(V)$，映射
\begin{equation}
 \Phi_Q:\End(V)\to\End(V),
 \qquad
 \Phi_Q(T)=Q^{-1}TQ
\end{equation}
是代数自同构。也就是说，
\begin{align}
 \Phi_Q(aS+bT)&=a\Phi_Q(S)+b\Phi_Q(T),\\
 \Phi_Q(ST)&=\Phi_Q(S)\Phi_Q(T),\\
 \Phi_Q(I)&=I.
\end{align}
其逆为 $\Phi_{Q^{-1}}$。
\end{proposition}

\begin{remark}
矩阵等价只保留秩，而矩阵相似保留线性变换的迭代结构，因而分类精细得多。相同特征多项式并不足以保证相似，例如
\begin{equation}
 \begin{pmatrix}0&1\\0&0\end{pmatrix}
 \quad\text{与}\quad
 \begin{pmatrix}0&0\\0&0\end{pmatrix}
\end{equation}
特征多项式同为 $\lambda^2$，但秩不同，所以不相似。
\end{remark}

\subsection{相似不变量}

若 $B=Q^{-1}AQ$，则
\begin{align}
 \rank B&=\rank A,\\
 \det B&=\det A,\\
 \tr B&=\tr A,\\
 \det(\lambda I-B)&=\det(\lambda I-A),\\
 f(B)&=Q^{-1}f(A)Q
\end{align}
对任意多项式 $f$ 成立。特别地，相似关系保持以下性质：
\begin{enumerate}
  \item 可逆性；
  \item 幂等性 $A^2=A$；
  \item 幂零性 $A^k=0$；
  \item 最小多项式与特征多项式；
  \item 各次幂的秩与核维数。
\end{enumerate}

\begin{remark}
相似一般不保持对称性、Hermite 性或正交性，因为这些性质还依赖于所选内积和基是否标准正交。若相似变换矩阵是正交或酉矩阵，则这些带内积的结构才会保持。
\end{remark}

\subsection{特征值与特征空间}

对 $A\in\F^{n\times n}$，若存在非零向量 $v$ 满足
\begin{equation}
 Av=\lambda v,
\end{equation}
则 $\lambda$ 称为特征值，$v$ 称为相应特征向量。特征空间为
\begin{equation}
 V_\lambda=\Ker(\lambda I-A).
\end{equation}
特征值由
\begin{equation}
 \chi_A(\lambda)=\det(\lambda I-A)=0
\end{equation}
确定。

\begin{theorem}[不同特征空间的直和性]
设 $\lambda_1,\ldots,\lambda_s$ 两两不同。若
\begin{equation}
 v_i\in V_{\lambda_i},
 \qquad i=1,\ldots,s,
\end{equation}
且每个 $v_i\neq0$，则 $v_1,\ldots,v_s$ 线性无关。更一般地，
\begin{equation}
 V_{\lambda_1}+\cdots+V_{\lambda_s}
\end{equation}
是直和。
\end{theorem}

\begin{proof}
设
\begin{equation}
 v_1+\cdots+v_s=0.
\end{equation}
对等式作用
\begin{equation}
 \prod_{j\neq i}(A-\lambda_jI),
\end{equation}
除第 $i$ 项外其余项都消失，得到
\begin{equation}
 \prod_{j\neq i}(\lambda_i-\lambda_j)v_i=0.
\end{equation}
因特征值两两不同，系数非零，故 $v_i=0$。
\end{proof}

\begin{corollary}
若 $A$ 在 $\F$ 中有 $n$ 个互不相同的特征值，则 $A$ 可对角化。
\end{corollary}

\begin{definition}
特征值 $\lambda$ 在特征多项式中的重数称为代数重数，记为 $m_a(\lambda)$；特征空间维数
\begin{equation}
 m_g(\lambda)=\dim\Ker(\lambda I-A)
\end{equation}
称为几何重数。
\end{definition}

\begin{theorem}
总有
\begin{equation}
 1\leq m_g(\lambda)\leq m_a(\lambda).
\end{equation}
矩阵 $A$ 可对角化当且仅当特征多项式在 $\F$ 中完全分裂，且每个特征值的几何重数等于代数重数。
\end{theorem}

\subsection{两个特征值计算例}

\begin{example}[可对角化矩阵及其幂]
原稿中的第一个矩阵为
\begin{equation}
 A_1=\begin{pmatrix}
 -1&3&-1\\
 -3&5&-1\\
 -3&3&1
 \end{pmatrix}.
\end{equation}
直接计算得
\begin{equation}
 \chi_{A_1}(\lambda)
 =\det(\lambda I-A_1)
 =(\lambda-2)^2(\lambda-1).
\end{equation}
当 $\lambda=2$ 时，可取两个线性无关特征向量
\begin{equation}
 v_1=(1,1,0)^T,
 \qquad
 v_2=(1,0,-3)^T.
\end{equation}
当 $\lambda=1$ 时，可取
\begin{equation}
 v_3=(1,1,1)^T.
\end{equation}
因此
\begin{equation}
 P=(v_1,v_2,v_3),
 \qquad
 A_1=P\diag(2,2,1)P^{-1},
\end{equation}
并且
\begin{equation}
 A_1^m=P\diag(2^m,2^m,1)P^{-1}.
\end{equation}
这比逐次相乘更适合计算高次幂。
\end{example}

\begin{example}[实矩阵的复特征值与实分块]
令
\begin{equation}
 A_2=\begin{pmatrix}
 4&-5&7\\
 1&-4&9\\
 -4&0&5
 \end{pmatrix}.
\end{equation}
计算得
\begin{equation}
 \chi_{A_2}(\lambda)
 =(\lambda-1)\bigl((\lambda-2)^2+3^2\bigr).
\end{equation}
所以在 $\R$ 上只有特征值 $1$，可取特征向量
\begin{equation}
 v_1=(1,2,1)^T,
\end{equation}
矩阵不能在 $\R$ 上对角化。在 $\C$ 上另有
\begin{equation}
 \lambda_2=2+3\ii,
 \qquad
 \lambda_3=2-3\ii,
\end{equation}
可取
\begin{equation}
 v_2=(3-3\ii,5-3\ii,4)^T,
 \qquad
 v_3=\overline{v_2}.
\end{equation}
于是 $A_2$ 在 $\C$ 上可对角化。

也可在实数域内令
\begin{equation}
 u=\Re v_2=(3,5,4)^T,
 \qquad
 w=\Im v_2=(-3,-3,0)^T.
\end{equation}
则 $\Span\{u,w\}$ 是实不变子空间，并且在基 $(u,w)$ 下限制矩阵为
\begin{equation}
 \begin{pmatrix}2&3\\-3&2\end{pmatrix}
\end{equation}
或按基次序不同得到其转置型。故 $A_2^m$ 可由一个一维块 $1^m$ 和该二维旋转--伸缩块的 $m$ 次幂计算。
\end{example}

\subsection{Gershgorin 圆盘定理}

\begin{theorem}[Gershgorin]
设 $A=(a_{ij})\in\C^{n\times n}$，并定义圆盘
\begin{equation}
 D_i=\set{z\in\C:
 |z-a_{ii}|\leq\sum_{j\neq i}|a_{ij}|}.
\end{equation}
则每个特征值都属于至少一个圆盘：
\begin{equation}
 \sigma(A)\subseteq\bigcup_{i=1}^nD_i.
\end{equation}
\end{theorem}

\begin{proof}
设 $Ax=\lambda x$，$x\neq0$，取 $k$ 使
\begin{equation}
 |x_k|=\max_i|x_i|.
\end{equation}
第 $k$ 行给出
\begin{equation}
 (\lambda-a_{kk})x_k
 =\sum_{j\neq k}a_{kj}x_j.
\end{equation}
因此
\begin{equation}
 |\lambda-a_{kk}|
 \leq\sum_{j\neq k}|a_{kj}|
 \frac{|x_j|}{|x_k|}
 \leq\sum_{j\neq k}|a_{kj}|,
\end{equation}
即 $\lambda\in D_k$。
\end{proof}

\begin{remark}
若若干 Gershgorin 圆盘与其余圆盘不相交，则该连通部分中包含的特征值个数等于其中圆盘的个数，按代数重数计算。这一加强结论可由连续变形到对角矩阵得到。
\end{remark}

\subsection{上三角化与 Schur 分解}

\begin{theorem}
若 $A\in\C^{n\times n}$，则存在可逆矩阵 $P$ 使
\begin{equation}
 P^{-1}AP
\end{equation}
为上三角矩阵。其对角元正是 $A$ 的特征值，按代数重数排列。
\end{theorem}

\begin{proof}
复数域上特征多项式有根。取一个特征向量作为新基的第一个向量，则矩阵具有分块形式
\begin{equation}
 \begin{pmatrix}\lambda&*\\0&A_1\end{pmatrix}.
\end{equation}
对 $A_1$ 归纳即可。
\end{proof}

\begin{theorem}[Schur 酉三角化]
对任意 $A\in\C^{n\times n}$，存在酉矩阵 $U$ 使
\begin{equation}
 U^*AU=T
\end{equation}
为上三角矩阵。
\end{theorem}

\begin{remark}
实矩阵未必有实特征值，因此未必能在 $\R$ 上三角化。但存在正交矩阵 $Q$，使 $Q^TAQ$ 为实拟上三角矩阵，其对角块为 $1\times1$ 实特征值块或对应一对共轭复特征值的 $2\times2$ 实块。这是实 Schur 分解。
\end{remark}

\section{习题课 \#11（11月22日）}

\subsection{正规矩阵与谱定理}

在复数域上，用 $A^*$ 表示共轭转置。

\begin{definition}
矩阵 $A\in\C^{n\times n}$ 称为正规矩阵，如果
\begin{equation}
 A^*A=AA^*.
\end{equation}
Hermite 矩阵、酉矩阵和复数域上的对角矩阵都是正规矩阵。
\end{definition}

\begin{theorem}[正规矩阵谱定理]
矩阵 $A\in\C^{n\times n}$ 正规，当且仅当存在酉矩阵 $U$ 与对角矩阵 $D$ 使
\begin{equation}
 A=UDU^*.
\end{equation}
\end{theorem}

\begin{proof}
若 $A=UDU^*$，则
\begin{equation}
 A^*A=U D^*D U^*=U DD^*U^*=AA^*.
\end{equation}

反之，由 Schur 分解，存在酉矩阵 $U$ 使
\begin{equation}
 T=U^*AU
\end{equation}
为上三角矩阵。正规性在酉相似下保持，所以 $T$ 正规。比较 $T^*T$ 与 $TT^*$ 的第一个对角元可知第一行非对角元素全为零，再对右下角子块归纳，得到 $T$ 必为对角矩阵。
\end{proof}

\begin{corollary}[Hermite 与实对称谱定理]
\begin{enumerate}
  \item 若 $A=A^*$，则 $A$ 的所有特征值均为实数，且可酉对角化；
  \item 若 $A\in\R^{n\times n}$ 且 $A=A^T$，则存在实正交矩阵 $Q$ 使
  \begin{equation}
   A=QDQ^T,
  \end{equation}
  其中 $D$ 为实对角矩阵。
\end{enumerate}
\end{corollary}

\begin{theorem}[不同特征值的特征向量正交]
若 $A$ 正规，且
\begin{equation}
 Ax=\lambda x,
 \qquad
 Ay=\mu y,
 \qquad
 \lambda\neq\mu,
\end{equation}
则
\begin{equation}
 x^*y=0.
\end{equation}
\end{theorem}

\begin{proof}
谱定理给出一组标准正交特征向量基，不同特征值对应的特征空间由互不相交的基向量子集张成，故彼此正交。也可利用正规矩阵满足
\begin{equation}
 \Ker(A-\lambda I)=\Ker(A^*-\overline\lambda I)
\end{equation}
直接计算。
\end{proof}

\subsection{实矩阵在实数域与复数域上的相似性}

\begin{theorem}
设 $A,B\in\R^{n\times n}$。则 $A$ 与 $B$ 在 $\R$ 上相似，当且仅当它们在 $\C$ 上相似。
\end{theorem}

\begin{proof}
实相似显然蕴含复相似。反之，假设存在复可逆矩阵
\begin{equation}
 S=P+\ii Q,
 \qquad P,Q\in\R^{n\times n},
\end{equation}
使
\begin{equation}
 AS=SB.
\end{equation}
比较实部与虚部得到
\begin{equation}
 AP=PB,
 \qquad
 AQ=QB.
\end{equation}
因此对每个 $t\in\R$，
\begin{equation}
 A(P+tQ)=(P+tQ)B.
\end{equation}
多项式 $p(t)=\det(P+tQ)$ 不恒为零，因为 $p(\ii)=\det(P+\ii Q)\neq0$。它在实数轴上只有有限多个零点，所以可选 $t\in\R$ 使 $P+tQ$ 可逆，从而得到实相似变换。
\end{proof}

\subsection{由一个矩阵生成的多项式代数}

\begin{definition}
对 $A\in\F^{n\times n}$，定义
\begin{equation}
 \F[A]=\set{f(A):f\in\F[t]},
\end{equation}
并定义 $A$ 的中心化子
\begin{equation}
 \mathcal C(A)=\set{B\in\F^{n\times n}:AB=BA}.
\end{equation}
\end{definition}

\begin{proposition}
$\F[A]$ 与 $\mathcal C(A)$ 都是 $\F^{n\times n}$ 的线性子空间，并且
\begin{equation}
 \F[A]\subseteq\mathcal C(A).
\end{equation}
若 $B\in\mathcal C(A)$，则
\begin{equation}
 \F[B]\subseteq\mathcal C(A).
\end{equation}
\end{proposition}

\begin{proof}
若 $AB=BA$，则归纳可得 $AB^k=B^kA$，继而 $A$ 与任意 $B$ 的多项式可交换。取 $B=A$ 即得第一项包含关系。
\end{proof}

\begin{proposition}[正定平方根是多项式]
设 $B\in\R^{n\times n}$ 为实对称正定矩阵。则存在多项式 $p$，使
\begin{equation}
 B^{1/2}=p(B),
\end{equation}
其中 $B^{1/2}$ 是唯一的实对称正定平方根。因此，若 $AB=BA$，则
\begin{equation}
 AB^{1/2}=B^{1/2}A.
\end{equation}
同理，$B^{-1/2}$ 也可写成 $B$ 的多项式。
\end{proposition}

\begin{proof}
设 $B$ 的互异特征值为 $\lambda_1,\ldots,\lambda_s>0$。由 Lagrange 插值，可取多项式 $p$ 满足
\begin{equation}
 p(\lambda_j)=\sqrt{\lambda_j},
 \qquad j=1,\ldots,s.
\end{equation}
若 $B=Q\diag(\lambda_1',\ldots,\lambda_n')Q^T$，则
\begin{equation}
 p(B)=Q\diag(\sqrt{\lambda_1'},\ldots,\sqrt{\lambda_n'})Q^T=B^{1/2}.
\end{equation}
\end{proof}

\begin{remark}
矩阵的三个常见分类关系分别为
\begin{equation}
 B=PAQ\quad\text{（等价）},
\end{equation}
\begin{equation}
 B=P^{-1}AP\quad\text{（相似）},
\end{equation}
以及
\begin{equation}
 B=P^TAP\quad\text{（合同）}.
\end{equation}
它们分别对应一般线性映射、线性变换和双线性型或二次型的换基。典型标准形分别由秩、Jordan 或有理标准形、以及惯性指数描述。
\end{remark}

\subsection{Fibonacci 数列的矩阵解法}

设
\begin{equation}
 a_0=0,
 \qquad a_1=1,
 \qquad a_{n+2}=a_{n+1}+a_n.
\end{equation}
令
\begin{equation}
 \alpha_n=\begin{pmatrix}a_{n+1}\\a_n\end{pmatrix},
 \qquad
 M=\begin{pmatrix}1&1\\1&0\end{pmatrix}.
\end{equation}
则
\begin{equation}
 \alpha_{n+1}=M\alpha_n,
 \qquad
 \alpha_n=M^n\alpha_0.
\end{equation}

矩阵 $M$ 的特征值为
\begin{equation}
 \varphi=\frac{1+\sqrt5}{2},
 \qquad
 \psi=\frac{1-\sqrt5}{2}=-\varphi^{-1}.
\end{equation}
对角化后得到 Binet 公式
\begin{equation}
 a_n=\frac{\varphi^n-\psi^n}{\sqrt5}.
\end{equation}
该方法同时说明递推数列、矩阵幂和特征值之间的直接联系。

\subsection{诱导矩阵二范数与谱半径}

\begin{definition}
对 $A\in\C^{m\times n}$，诱导二范数定义为
\begin{equation}
 \norm{A}_2
 =\sup_{x\neq0}\frac{\norm{Ax}_2}{\norm x_2}.
\end{equation}
对方阵 $A$，谱半径定义为
\begin{equation}
 \rho(A)=\max_{\lambda\in\sigma(A)}|\lambda|.
\end{equation}
\end{definition}

\begin{proposition}
诱导二范数满足：
\begin{enumerate}
  \item $0\leq\norm A_2<\infty$，且 $\norm A_2=0$ 当且仅当 $A=0$；
  \item $\norm{aA}_2=|a|\norm A_2$；
  \item $\norm{A+B}_2\leq\norm A_2+\norm B_2$；
  \item $\norm{AB}_2\leq\norm A_2\norm B_2$；
  \item $\norm A_2^2=\rho(A^*A)$。
\end{enumerate}
\end{proposition}

\begin{theorem}
对任意方阵，
\begin{equation}
 \rho(A)\leq\norm A_2.
\end{equation}
若 $A$ 正规，则
\begin{equation}
 \rho(A)=\norm A_2.
\end{equation}
\end{theorem}

\begin{proof}
若 $Av=\lambda v$，则
\begin{equation}
 |\lambda|\norm v_2=\norm{Av}_2\leq\norm A_2\norm v_2.
\end{equation}
对正规矩阵写成 $A=UDU^*$，酉不变性给出
\begin{equation}
 \norm A_2=\norm D_2=\max_i|d_i|=\rho(A).
\end{equation}
\end{proof}

\begin{remark}
一般矩阵中不一定取等号。例如
\begin{equation}
 N=\begin{pmatrix}0&1\\0&0\end{pmatrix}
\end{equation}
满足 $\rho(N)=0$，但 $\norm N_2=1$。
\end{remark}

\subsection{矩阵指数与谱半径为零的判别}

由次乘法性，级数
\begin{equation}
 \e^A=\sum_{k=0}^{\infty}\frac{A^k}{k!}
\end{equation}
绝对收敛，因为
\begin{equation}
 \sum_{k=0}^{\infty}\frac{\norm A_2^k}{k!}
 \leq\e^{\norm A_2}.
\end{equation}
因此矩阵指数定义良好，并满足 $AB=BA$ 时
\begin{equation}
 \e^{A+B}=\e^A\e^B.
\end{equation}

\begin{theorem}
对 $A\in\C^{n\times n}$，下列条件等价：
\begin{enumerate}
  \item $\rho(A)=0$；
  \item $A$ 的全部特征值为 $0$；
  \item $A$ 幂零；
  \item 对每个 $k\geq1$，$\tr(A^k)=0$。
\end{enumerate}
\end{theorem}

\begin{proof}
前三项由 Cayley--Hamilton 定理直接等价。若特征值全为零，则 $\tr(A^k)$ 是特征值 $k$ 次幂之和，故为零。

反之，设互异非零特征值为 $\lambda_1,\ldots,\lambda_s$，其代数重数为 $m_1,\ldots,m_s$。由
\begin{equation}
 \sum_{j=1}^s m_j\lambda_j^k=0,
 \qquad k=1,\ldots,s,
\end{equation}
除去一个 $\lambda_j$ 后得到 Vandermonde 型线性方程组。其系数矩阵可逆，迫使 $m_j\lambda_j=0$，矛盾。因此不存在非零特征值。
\end{proof}

\subsection{半正定矩阵的等价刻画}

\begin{theorem}
对实方阵 $A$，下列条件等价：
\begin{enumerate}
  \item 存在实矩阵 $B$ 使
  \begin{equation}
   A=B^TB;
  \end{equation}
  \item $A=A^T$ 且
  \begin{equation}
   x^TAx\geq0
   \qquad\text{对所有 }x\in\R^n;
  \end{equation}
  \item 存在唯一的实对称半正定矩阵 $C$ 使
  \begin{equation}
   A=C^2.
  \end{equation}
\end{enumerate}
\end{theorem}

\begin{proof}
若 $A=B^TB$，则 $A$ 对称且
\begin{equation}
 x^TAx=\norm{Bx}_2^2\geq0.
\end{equation}
若 $A$ 对称半正定，谱定理给出
\begin{equation}
 A=Q\diag(\lambda_1,\ldots,\lambda_n)Q^T,
 \qquad \lambda_i\geq0.
\end{equation}
取
\begin{equation}
 C=Q\diag(\sqrt{\lambda_1},\ldots,\sqrt{\lambda_n})Q^T
\end{equation}
即可。半正定平方根的唯一性可在每个特征空间上验证。最后令 $B=C$ 即回到第一项。
\end{proof}

\section{习题课 \#12（11月29日）}

\subsection{实斜对称矩阵的合同标准形}

令
\begin{equation}
 J=\begin{pmatrix}0&1\\-1&0\end{pmatrix}.
\end{equation}

\begin{theorem}
设 $A\in\F^{n\times n}$ 满足 $A^T=-A$，且 $\operatorname{char}\F\neq2$。若 $\rank A=2r$，则存在可逆矩阵 $P$ 使
\begin{equation}
 P^TAP=\diag(\underbrace{J,\ldots,J}_{r\text{ 个}},0_{n-2r}).
\end{equation}
特别地，斜对称矩阵的秩为偶数。
\end{theorem}

\begin{proof}
对 $n$ 归纳。若 $A=0$ 则结论显然。否则存在 $i\neq j$ 使 $a_{ij}\neq0$。同时置换第 $i,j$ 行和列，并作一个对角缩放，可使左上角 $2\times2$ 块变为 $J$。此时
\begin{equation}
 A=\begin{pmatrix}J&B\\-B^T&C\end{pmatrix},
 \qquad C^T=-C.
\end{equation}
利用保持合同关系的分块消元，可消去 $B$，得到
\begin{equation}
 A\sim_{\mathrm{cong}}\begin{pmatrix}J&0\\0&C+B^TJB\end{pmatrix}.
\end{equation}
右下角仍为斜对称矩阵，对它应用归纳假设即可。
\end{proof}

\subsection{二次型分解为两个一次型}

设实二次型
\begin{equation}
 q(x)=x^TAx,
 \qquad A=A^T\in\R^{n\times n}.
\end{equation}

\begin{theorem}
非零二次型 $q$ 可以分解为两个实一次齐次多项式之积
\begin{equation}
 q(x)=(\alpha^Tx)(\beta^Tx)
\end{equation}
当且仅当下列两种情形之一成立：
\begin{enumerate}
  \item $\rank A=1$；
  \item $\rank A=2$，且 $A$ 的正、负惯性指数均为 $1$。
\end{enumerate}
零二次型当然也可作平凡分解。
\end{theorem}

\begin{proof}
若有分解，则
\begin{equation}
 A=\frac12(\alpha\beta^T+\beta\alpha^T),
\end{equation}
所以 $\rank A\leq2$。若 $\alpha,\beta$ 线性相关，则 $q$ 是一个一次型平方的常数倍，秩为 $1$。若二者线性无关，则在它们张成的二维空间上可换基写成
\begin{equation}
 q=y_1y_2
 =\frac14\bigl((y_1+y_2)^2-(y_1-y_2)^2\bigr),
\end{equation}
故正、负惯性指数各为 $1$。

反之，秩 $1$ 时由 Sylvester 惯性定理可写成 $q=\pm y_1^2$，显然可分解。秩 $2$ 且惯性为 $(1,1)$ 时可写成
\begin{equation}
 q=y_1^2-y_2^2=(y_1+y_2)(y_1-y_2).
\end{equation}
\end{proof}

\subsection{Rayleigh 商与极值特征值}

设 $A=A^T\in\R^{n\times n}$，特征值按
\begin{equation}
 \lambda_1\geq\lambda_2\geq\cdots\geq\lambda_n
\end{equation}
排列。对 $x\neq0$ 定义 Rayleigh 商
\begin{equation}
 R_A(x)=\frac{x^TAx}{x^Tx}.
\end{equation}

\begin{theorem}
有
\begin{equation}
 \max_{x\neq0}R_A(x)=\lambda_1,
 \qquad
 \min_{x\neq0}R_A(x)=\lambda_n.
\end{equation}
等号分别在最大、最小特征值对应的特征空间上取得。
\end{theorem}

\begin{proof}
由谱定理，$A=QDQ^T$，其中 $D=\diag(\lambda_1,\ldots,\lambda_n)$。令 $y=Q^Tx$，则
\begin{equation}
 R_A(x)=\frac{\sum_{i=1}^n\lambda_i y_i^2}
 {\sum_{i=1}^n y_i^2},
\end{equation}
它是各特征值的加权平均，故位于 $[\lambda_n,\lambda_1]$ 内。
\end{proof}

\subsection{Courant--Fischer 极小极大原理}

\begin{theorem}[Courant--Fischer]
对 $k=1,\ldots,n$，有
\begin{equation}
 \lambda_k
 =\max_{\substack{S\leq\R^n\\\dim S=k}}
 \ \min_{\substack{x\in S\\x\neq0}}R_A(x)
\end{equation}
以及
\begin{equation}
 \lambda_k
 =\min_{\substack{T\leq\R^n\\\dim T=n-k+1}}
 \ \max_{\substack{x\in T\\x\neq0}}R_A(x).
\end{equation}
\end{theorem}

\begin{proof}
设 $v_1,\ldots,v_n$ 是按上述顺序排列的标准正交特征向量。取
\begin{equation}
 S_0=\Span(v_1,\ldots,v_k),
\end{equation}
则 $S_0$ 上 Rayleigh 商的最小值为 $\lambda_k$。另一方面，任意 $k$ 维子空间 $S$ 与
\begin{equation}
 \Span(v_k,\ldots,v_n)
\end{equation}
维数之和大于 $n$，故有非零交点。在该交点上 $R_A(x)\leq\lambda_k$。这证明第一个公式。第二个公式同理，或对 $-A$ 应用第一个公式。
\end{proof}

\subsection{可交换矩阵的同时标准化}

\begin{theorem}[同时上三角化]
若 $A_1,\ldots,A_s\in\C^{n\times n}$ 两两可交换，则存在同一个可逆矩阵 $P$，使
\begin{equation}
 P^{-1}A_jP
\end{equation}
对所有 $j$ 都是上三角矩阵。
\end{theorem}

\begin{proof}
先证明存在公共特征向量。取 $A_1$ 的一个特征空间 $E$。由于 $A_jA_1=A_1A_j$，每个 $A_j$ 都保持 $E$ 不变。在 $E$ 上继续取 $A_2$ 的特征空间，如此反复，可得到非零公共不变子空间，最终取得公共特征向量。以该向量作为第一基向量后，对商空间归纳。
\end{proof}

\begin{theorem}[同时对角化]
若 $A_1,\ldots,A_s\in\C^{n\times n}$ 两两可交换且各自可对角化，则它们可被同一个可逆矩阵同时对角化。
\end{theorem}

\begin{proof}
把空间分解为 $A_1$ 的特征空间直和。可交换性保证每个特征空间都被其余矩阵保持。其余矩阵在这些不变子空间上的限制仍可对角化，对每个子空间递归处理并合并所选基即可。
\end{proof}

\begin{corollary}
若 $A_1,\ldots,A_s$ 是两两可交换的实对称矩阵，则存在同一个正交矩阵 $Q$ 使所有
\begin{equation}
 Q^TA_jQ
\end{equation}
均为实对角矩阵。
\end{corollary}

\subsection{平方和表示与惯性指数的估计}

设二次型写成
\begin{equation}
 q(x)=\sum_{j=1}^s(\alpha_j^Tx)^2
 -\sum_{j=s+1}^{s+t}(\alpha_j^Tx)^2.
\end{equation}

\begin{proposition}
$q$ 的正惯性指数不超过 $s$，负惯性指数不超过 $t$。
\end{proposition}

\begin{proof}
令 $B$ 的各行为 $\alpha_1^T,\ldots,\alpha_s^T$，则
\begin{equation}
 \dim\Ker B\geq n-s.
\end{equation}
在 $\Ker B$ 上，前 $s$ 个正平方项全为零，所以 $q\leq0$。若正惯性指数大于 $s$，则存在维数大于 $s$ 的子空间 $V$，使 $q$ 在 $V\setminus\{0\}$ 上严格为正。由维数公式，$V\cap\Ker B$ 含非零向量，矛盾。负惯性指数的估计对 $-q$ 使用同一论证。
\end{proof}

\subsection{奇异值分解与极分解}

\begin{theorem}[奇异值分解]
对任意 $A\in\C^{m\times n}$，存在酉矩阵
\begin{equation}
 U\in\C^{m\times m},
 \qquad V\in\C^{n\times n}
\end{equation}
以及矩形对角矩阵
\begin{equation}
 \Sigma=\diag(\sigma_1,\ldots,\sigma_r,0,\ldots),
 \qquad
 \sigma_1\geq\cdots\geq\sigma_r>0,
\end{equation}
使
\begin{equation}
 A=U\Sigma V^*.
\end{equation}
非零 $\sigma_i$ 是 $A^*A$ 的正特征值平方根，称为奇异值。
\end{theorem}

\begin{proof}
对正半定矩阵 $A^*A$ 作酉对角化：
\begin{equation}
 V^*A^*AV=\diag(\sigma_1^2,\ldots,\sigma_r^2,0,\ldots,0).
\end{equation}
令 $v_i$ 为 $V$ 的列，并取
\begin{equation}
 u_i=\sigma_i^{-1}Av_i,
 \qquad i=1,\ldots,r.
\end{equation}
这些向量标准正交，把它们补成 $\C^m$ 的标准正交基即可得到 $U$。
\end{proof}

\begin{theorem}[极分解]
任意 $A\in\C^{m\times n}$ 可写成
\begin{equation}
 A=QH,
 \qquad
 H=(A^*A)^{1/2}\succeq0,
\end{equation}
其中 $Q$ 在 $\Image H$ 上是等距映射，并可延拓为适当的部分酉矩阵。若 $A$ 为可逆方阵，则 $Q$ 为酉矩阵，而且分解唯一：
\begin{equation}
 Q=A(A^*A)^{-1/2}.
\end{equation}
\end{theorem}

\begin{proof}
由 SVD $A=U\Sigma V^*$，取
\begin{equation}
 H=V\Sigma V^*,
 \qquad
 Q=UV^*
\end{equation}
在方阵可逆情形即得结论。矩形或退化情形中只在非零奇异子空间上定义对应的部分等距映射，再作延拓。
\end{proof}

\subsection{二次型的配方法与合同对角化}

设
\begin{equation}
 q(x)=x^TAx,
 \qquad A=A^T,
 \qquad \operatorname{char}\F\neq2.
\end{equation}
若 $a_{11}\neq0$，则
\begin{equation}
 q(x)=a_{11}
 \left(x_1+\sum_{j=2}^n\frac{a_{1j}}{a_{11}}x_j\right)^2
 +q_1(x_2,\ldots,x_n),
\end{equation}
从而把维数降低一阶。若所有对角元均为零但某个 $a_{ij}\neq0$，可先作变量替换
\begin{equation}
 y_i=x_i+x_j,
 \qquad
 y_j=x_i-x_j,
\end{equation}
使新表达式出现非零平方项，再继续配方。由此可逐步得到合同对角形。

\begin{example}
将
\begin{equation}
 q(x_1,x_2,x_3)=x_1x_2+x_1x_3-3x_2x_3
\end{equation}
化为标准形。先令
\begin{equation}
 x_1=y_1+y_2,
 \qquad
 x_2=y_1-y_2,
 \qquad
 x_3=y_3.
\end{equation}
则
\begin{align}
 q&=y_1^2-y_2^2-2y_1y_3+4y_2y_3\\
  &=(y_1-y_3)^2-(y_2-2y_3)^2+3y_3^2.
\end{align}
再令
\begin{equation}
 z_1=y_1-y_3,
 \qquad
 z_2=y_2-2y_3,
 \qquad
 z_3=y_3,
\end{equation}
得到
\begin{equation}
 q=z_1^2-z_2^2+3z_3^2.
\end{equation}
相应的可逆变量关系为
\begin{equation}
 x_1=z_1+z_2+3z_3,
 \qquad
 x_2=z_1-z_2-z_3,
 \qquad
 x_3=z_3.
\end{equation}
\end{example}

\subsection{对偶空间}

\begin{definition}
向量空间 $V$ 的对偶空间定义为
\begin{equation}
 V^*=\Hom(V,\F),
\end{equation}
即 $V$ 上所有线性函数的空间。
\end{definition}

若 $V$ 有基 $e_1,\ldots,e_n$，定义线性函数 $e_1^*,\ldots,e_n^*$ 使
\begin{equation}
 e_i^*(e_j)=\delta_{ij}.
\end{equation}
则它们构成 $V^*$ 的一组基，称为对偶基。因此有限维时
\begin{equation}
 \dim V^*=\dim V,
\end{equation}
并且一旦选择基便得到一个同构 $V\cong V^*$。不过该同构依赖于基或内积，并非自然给定。

相反，映射
\begin{equation}
 \iota:V\longrightarrow V^{**},
 \qquad
 \iota(v)(f)=f(v)
\end{equation}
不依赖任何基。有限维时它是同构，称为典范同构。

\begin{example}[无限维时 $V$ 与 $V^*$ 的巨大差别]
令
\begin{equation}
 V=\R^{(\N)}
 =\set{(a_i)_{i\geq0}:a_i\text{ 只有有限个非零项}}.
\end{equation}
它有可数 Hamel 基 $e_0,e_1,\ldots$。任意序列 $b=(b_i)_{i\geq0}$ 都定义一个线性函数
\begin{equation}
 f_b(a)=\sum_{i\geq0}b_i a_i,
\end{equation}
因为对每个 $a\in V$，该和只有有限项。反之，每个线性函数都由 $f(e_i)$ 唯一确定，所以
\begin{equation}
 V^*\cong\R^{\N},
\end{equation}
即所有实数序列的空间。

而 $\R^{\N}$ 不存在可数 Hamel 基。事实上，对不同实数 $t$，序列
\begin{equation}
 v(t)=(1,t,t^2,t^3,\ldots)
\end{equation}
构成一个不可数线性无关族，因为任取有限多个不同的 $t$，前若干坐标形成 Vandermonde 矩阵。故在无限维情形，$V$ 与 $V^*$ 通常不再同构。
\end{example}

\section{习题课 \#13（12月6日）}

\subsection{正定矩阵的判别}

设 $A=A^T\in\R^{n\times n}$。下列条件等价：
\begin{enumerate}
  \item 对每个 $x\neq0$，
  \begin{equation}
   x^TAx>0;
  \end{equation}
  \item $A$ 的全部特征值均为正；
  \item 存在可逆矩阵 $P$ 使
  \begin{equation}
   A=P^TP;
  \end{equation}
  \item 存在对角元为正的可逆下三角矩阵 $L$ 使
  \begin{equation}
   A=LL^T;
  \end{equation}
  \item 所有顺序主子式均为正；
  \item 所有主子式均为正。
\end{enumerate}
其中第四项是 Cholesky 分解，第五项通常称为 Sylvester 判别法。

\begin{proof}
由实对称谱定理，
\begin{equation}
 A=Q\diag(\lambda_1,\ldots,\lambda_n)Q^T.
\end{equation}
于是
\begin{equation}
 x^TAx=\sum_i\lambda_i y_i^2,
 \qquad y=Q^Tx,
\end{equation}
所以前两项等价。若所有 $\lambda_i>0$，取
\begin{equation}
 P=\diag(\sqrt{\lambda_1},\ldots,\sqrt{\lambda_n})Q^T
\end{equation}
便有 $A=P^TP$；反向由 $x^TP^TPx=\norm{Px}^2>0$ 得到。

Cholesky 分解可按 Schur 补归纳构造。写
\begin{equation}
 A=\begin{pmatrix}a&b^T\\b&C\end{pmatrix},
 \qquad a>0.
\end{equation}
则 $A$ 正定当且仅当
\begin{equation}
 C-a^{-1}bb^T
\end{equation}
正定。对该 Schur 补递归分解即可。分解中前 $k$ 阶顺序主子式等于前 $k$ 个 Cholesky 对角元平方的乘积，因此正定蕴含所有顺序主子式为正；反向可按同一递推构造 Cholesky 分解。所有主子式为正显然蕴含顺序主子式为正，而正定矩阵的任意主子矩阵仍正定，故其主子式均正。
\end{proof}

\subsection{两个正定性例题}

\begin{example}
考虑
\begin{equation}
 f(x_1,x_2,x_3)
 =99x_1^2-12x_1x_2+48x_1x_3
 +130x_2^2-60x_2x_3+71x_3^2.
\end{equation}
其对称矩阵为
\begin{equation}
 A=\begin{pmatrix}
 99&-6&24\\
 -6&130&-30\\
 24&-30&71
 \end{pmatrix}.
\end{equation}
顺序主子式为
\begin{equation}
 \Delta_1=99>0,
 \qquad
 \Delta_2=99\cdot130-36=12834>0,
\end{equation}
以及
\begin{equation}
 \Delta_3=\det A=755874>0.
\end{equation}
故 $A$ 正定，二次型 $f$ 为正定二次型。
\end{example}

\begin{example}
考虑参数矩阵
\begin{equation}
 B(t)=\begin{pmatrix}
 1&t&5\\
 t&4&3\\
 5&3&1
 \end{pmatrix}.
\end{equation}
若 $B(t)$ 正定，则前两个顺序主子式要求
\begin{equation}
 4-t^2>0,
 \qquad -2<t<2.
\end{equation}
但
\begin{equation}
 \det B(t)=-t^2+30t-105
 =-(t-15)^2+120.
\end{equation}
其为正要求
\begin{equation}
 15-2\sqrt{30}<t<15+2\sqrt{30},
\end{equation}
与 $(-2,2)$ 没有交集。因此不存在实数 $t$ 使 $B(t)$ 正定。
\end{example}

\subsection{Cholesky 分解}

\begin{theorem}[Cholesky]
若 $A=A^T\in\R^{n\times n}$ 正定，则存在唯一的下三角矩阵 $L$，其对角元全正，并满足
\begin{equation}
 A=LL^T.
\end{equation}
\end{theorem}

\begin{proof}
存在性已可由上一节的 Schur 补递推得到。为说明唯一性，若
\begin{equation}
 A=L_1L_1^T=L_2L_2^T,
\end{equation}
则
\begin{equation}
 L_2^{-1}L_1=L_2^TL_1^{-T}.
\end{equation}
左端是对角元全正的下三角矩阵，右端是上三角矩阵，所以二者必须是对角矩阵。再由
\begin{equation}
 (L_2^{-1}L_1)(L_2^{-1}L_1)^T=I
\end{equation}
及正对角元可知该对角矩阵为 $I$，故 $L_1=L_2$。
\end{proof}

\subsection{利用迹刻画正定性}

\begin{theorem}
设 $A=A^T\in\R^{n\times n}$ 可逆。则 $A$ 正定，当且仅当对每个正定矩阵 $B$ 都有
\begin{equation}
 \tr(AB)>0.
\end{equation}
\end{theorem}

\begin{proof}
若 $A$ 正定，写成
\begin{equation}
 A=P^TP,
 \qquad B=Q^TQ,
\end{equation}
其中 $P,Q$ 可逆。则
\begin{equation}
 \tr(AB)=\tr(P^TPQ^TQ)
 =\tr\bigl((QP^T)(QP^T)^T\bigr)
 =\norm{QP^T}_F^2>0.
\end{equation}

反之，若存在 $x\neq0$ 使 $x^TAx<0$，取
\begin{equation}
 B=xx^T+\varepsilon I
\end{equation}
并令 $\varepsilon>0$ 足够小，则 $B$ 正定而
\begin{equation}
 \tr(AB)=x^TAx+\varepsilon\tr A<0,
\end{equation}
矛盾。因此 $A$ 半正定。结合 $A$ 可逆可知 $A$ 正定。
\end{proof}

\subsection{实对称矩阵的相似与正交相似}

\begin{theorem}
设 $A,B\in\R^{n\times n}$ 都是实对称矩阵。则
\begin{equation}
 A\text{ 与 }B\text{ 相似}
 \quad\Longleftrightarrow\quad
 A\text{ 与 }B\text{ 正交相似}.
\end{equation}
\end{theorem}

\begin{proof}
正交相似当然蕴含相似。反之，设
\begin{equation}
 AP=PB
\end{equation}
且 $P$ 可逆。转置后有
\begin{equation}
 P^TA=BP^T.
\end{equation}
因此
\begin{equation}
 B(P^TP)=P^TAP=(P^TP)B.
\end{equation}
令
\begin{equation}
 H=(P^TP)^{1/2}.
\end{equation}
由多项式插值或函数演算，$H$ 也与 $B$ 可交换。取极分解
\begin{equation}
 P=QH,
 \qquad Q\text{ 正交}.
\end{equation}
代入 $AP=PB$ 得
\begin{equation}
 AQH=QHB=QBH.
\end{equation}
右乘 $H^{-1}$，得到 $AQ=QB$，即
\begin{equation}
 A=QBQ^T.
\end{equation}
\end{proof}

\subsection{秩一修正下的惯性变化}

对实对称矩阵 $M$，记
\begin{equation}
 n_+(M),\quad n_-(M),\quad n_0(M)
\end{equation}
分别为正、负、零特征值个数，并定义符号差
\begin{equation}
 s(M)=n_+(M)-n_-(M).
\end{equation}

\begin{theorem}
设 $A=A^T\in\R^{n\times n}$ 可逆，$\alpha\in\R^n$，并令
\begin{equation}
 B=A-\alpha\alpha^T,
 \qquad
 r=\alpha^TA^{-1}\alpha.
\end{equation}
则
\begin{equation}
 s(A)=s(B)+1-\sgn(1-r),
\end{equation}
其中约定 $\sgn(0)=0$。也就是说，
\begin{equation}
 s(A)=
 \begin{cases}
 s(B),&r<1,\\
 s(B)+1,&r=1,\\
 s(B)+2,&r>1.
 \end{cases}
\end{equation}
当 $r=1$ 时，$B$ 必奇异。
\end{theorem}

\begin{proof}
考虑对称分块矩阵
\begin{equation}
 M=\begin{pmatrix}1&\alpha^T\\\alpha&A\end{pmatrix}.
\end{equation}
以左上角 $1$ 作 Schur 补，得到
\begin{equation}
 M\sim_{\mathrm{cong}}\diag(1,B),
\end{equation}
故
\begin{equation}
 s(M)=1+s(B).
\end{equation}
另一方面，以 $A$ 作 Schur 补，得到
\begin{equation}
 M\sim_{\mathrm{cong}}\diag(1-r,A),
\end{equation}
所以
\begin{equation}
 s(M)=\sgn(1-r)+s(A).
\end{equation}
比较两式即可。
\end{proof}

\subsection{秩一全一修正矩阵的惯性}

令
\begin{equation}
 A=\lambda I_n+\ones\ones^T.
\end{equation}
向量 $\ones$ 是特征向量，对应特征值 $\lambda+n$；在 $\ones^\perp$ 上，$A$ 等于 $\lambda I$。因此特征值为
\begin{equation}
 \lambda+n\quad\text{一次},
 \qquad
 \lambda\quad\text{$n-1$ 次}.
\end{equation}
符号差为
\begin{equation}
 s(A)=
 \begin{cases}
 n,&\lambda>0,\\
 1,&\lambda=0,\\
 2-n,&-n<\lambda<0,\\
 1-n,&\lambda=-n,\\
 -n,&\lambda<-n.
 \end{cases}
\end{equation}
这也给出正、负惯性指数及奇异参数值。

\subsection{两个特殊二次型}

\begin{example}
令
\begin{equation}
 q(x)=\sum_{1\leq j<k\leq n}(-1)^{j+k}x_jx_k.
\end{equation}
取
\begin{equation}
 \alpha=(1,-1,1,-1,\ldots)^T.
\end{equation}
则 $q(x)=x^TAx$ 的矩阵满足
\begin{equation}
 A=-\frac12I+\frac12\alpha\alpha^T.
\end{equation}
因为 $\alpha^T\alpha=n$，$A$ 在 $\alpha$ 方向上的特征值为
\begin{equation}
 \frac{n-1}{2},
\end{equation}
在 $\alpha^\perp$ 上的特征值为 $-1/2$。因此 $n>1$ 时其惯性为
\begin{equation}
 (n_+,n_-,n_0)=(1,n-1,0).
\end{equation}
\end{example}

\begin{example}
令
\begin{equation}
 g(x)=\sum_{1\leq j<k\leq n}|j-k|x_jx_k.
\end{equation}
其对称矩阵 $D$ 的对角元为 $0$，非对角元为
\begin{equation}
 D_{jk}=\frac12|j-k|.
\end{equation}
依次对相邻行、列作差，可将 $D$ 合同化为一个正的一维块与 $n-1$ 个负的一维块。因此
\begin{equation}
 (n_+(D),n_-(D),n_0(D))=(1,n-1,0),
 \qquad n\geq2.
\end{equation}
原稿中的连续合同行列变换正是这一结论的直接计算证明。
\end{example}


\end{document}
